% Luc Destombe --- licence CC BY-NC-SA 4.0
% https://creativecommons.org/licenses/by-nc-sa/4.0/deed.fr
% Fichier autonome : compiler avec pdflatex, deux fois (signets, nombre de pages).
\documentclass[11pt,a4paper]{article}
\makeatletter
\newif\iflycee@unecolonne
\lycee@unecolonnetrue
\newif\iflycee@palatino
\lycee@palatinotrue

\RequirePackage[utf8]{inputenc}
\RequirePackage[T1]{fontenc}
\RequirePackage[french]{babel}
\RequirePackage{etoolbox}

\newcommand{\ShorthandsOff}{\@ifundefined{shorthandoff}{}{\shorthandoff{;:!?}}}
\newcommand{\ShorthandsOn}{\@ifundefined{shorthandon}{}{\shorthandon{;:!?}}}
\ShorthandsOff
\AtBeginDocument{\ShorthandsOn}
\AtBeginEnvironment{tikzpicture}{\ShorthandsOff}

\RequirePackage{geometry}
\RequirePackage{fancyhdr}
\RequirePackage{lastpage}
\RequirePackage{multicol}
\RequirePackage{array}
\RequirePackage{enumitem}
\RequirePackage{xcolor}
\RequirePackage{needspace}

\RequirePackage{amsmath,amssymb}
\RequirePackage{mathpazo}
\DeclareMathSizes{10.95}{10.95}{8}{5.5}
\begingroup\catcode`\_=\active \catcode`\^=\active
\gdef_#1{\sb{\scriptscriptstyle #1}}
\gdef^#1{\sp{\scriptscriptstyle\lycee@moinsepais #1}}
\endgroup
\newcommand\lycee@moins{\mathbin{%
  \setbox\z@\hbox{$\m@th\scriptscriptstyle\lycee@moinsorig$}%
  \dimen@=.55\wd\z@ \advance\dimen@-.5pt
  \hbox to.75\wd\z@{\hss$\m@th\scriptscriptstyle\vcenter{\hbox{%
    \tikz\draw[line width=.5pt,line cap=round](0,0)--(\the\dimen@,0);}}$\hss}}}
\begingroup\lccode`\~=`\- \lowercase{\endgroup\let~}\lycee@moins
\newcommand\lycee@moinsepais{\mathcode`\-="8000 }
\AtBeginDocument{\mathchardef\lycee@moinsorig=\mathcode`\-\relax}
\AtBeginDocument{\catcode`\_=12 \mathcode`\_="8000
  \catcode`\^=12 \mathcode`\^="8000 }
\newcommand\lycee@exposants{%
  \fontdimen13\textfont2=.48\fontdimen6\textfont2
  \fontdimen14\textfont2=.48\fontdimen6\textfont2
  \fontdimen15\textfont2=.48\fontdimen6\textfont2
  \fontdimen13\scriptfont2=.48\fontdimen6\scriptfont2
  \fontdimen14\scriptfont2=.48\fontdimen6\scriptfont2
  \fontdimen15\scriptfont2=.48\fontdimen6\scriptfont2}
\every@math@size\expandafter{\the\every@math@size\lycee@exposants}
\newlength{\retraitcalculs}
\setlength{\retraitcalculs}{2em}

\RequirePackage{tikz}
\usetikzlibrary{arrows.meta,calc,babel,shapes.misc}
\RequirePackage[most]{tcolorbox}
\tcbuselibrary{skins,breakable}

\definecolor{coulDef}{HTML}{1F4E79}
\definecolor{coulProp}{HTML}{7030A0}
\definecolor{coulRem}{HTML}{C55A11}
\definecolor{coulEx}{HTML}{2E7D32}
\definecolor{coulExo}{HTML}{B03A2E}
\definecolor{coulPart}{HTML}{1F4E79}
\definecolor{coulMeth}{HTML}{1B6E4F}
\definecolor{coulCourbe}{HTML}{0B5ED7}
\definecolor{coulCourbeB}{HTML}{C00000}
\definecolor{coulCourbeC}{HTML}{2E7D32}
\definecolor{coulGrille}{HTML}{8C8C8C}
\definecolor{coulRep}{HTML}{1B6E4F}
\definecolor{coulBar}{HTML}{2E7D32}

\newcommand{\R}{\mathbb{R}}
\newcommand{\Rs}{\mathbb{R}^{*}}
\renewcommand{\le}{\leqslant}
\renewcommand{\ge}{\geqslant}
\renewcommand{\approx}{\simeq}
\newcommand{\vect}[1]{\overrightarrow{#1}}
\newcommand{\norme}[1]{\left\lVert #1 \right\rVert}
\newcommand{\intervalle}[2]{\left[#1\,;#2\right]}
\RequirePackage{cancel}
\renewcommand{\CancelColor}{\color{coulExo}}
\newcommand{\fois}[1]{\times{\color{coulExo}#1}}
\newcommand{\barre}[1]{\cancel{#1}}

\tikzset{grille/.style={coulGrille,line width=0.4pt}}
\newcommand{\quadrillage}[4]{\draw[grille,step=1] (#1,#2) grid (#3,#4);}
\tikzset{
  croix/.style={cross out,draw,line width=0.6pt,inner sep=0pt,minimum size=4pt},
  croixdroite/.style={croix,rotate=45,line width=0.8pt},
}
\newcommand{\pt}[3]{\path (#1) node[croix]{} node[#2,font=\small,inner sep=2.5pt]{$#3$};}

\newcommand{\lycee@repere}[4]{%
  \draw[grille,step=1] (#1,#2) grid (#3,#4);
  \draw[-{Stealth[length=2mm]},thick] (#1,0)--({#3+0.4},0) node[below left=-1pt]{$x$};
  \draw[-{Stealth[length=2mm]},thick] (0,#2)--(0,{#4+0.4}) node[below left=-1pt]{$y$};
  \node[below left,font=\scriptsize,inner sep=1.5pt] at (0,0) {$0$};}
\newcommand{\lycee@gradx}[1]{\foreach \k/\l in {#1}{%
  \draw (\k,0.07)--(\k,-0.07) node[below,font=\scriptsize,inner sep=1.5pt]{$\l$};}}
\newcommand{\lycee@grady}[1]{\foreach \k/\l in {#1}{%
  \draw (0.07,\k)--(-0.07,\k) node[left,font=\scriptsize,inner sep=1.5pt]{$\l$};}}
\newcommand{\lycee@intff}[2]{\left[#1\,;#2\right]}
\newcommand{\lycee@intfo}[2]{\left[#1\,;#2\right[}
\newcommand{\lycee@intof}[2]{\left]#1\,;#2\right]}
\newcommand{\lycee@intoo}[2]{\left]#1\,;#2\right[}
\newcommand{\lycee@ens}[1]{\left\{#1\right\}}
\AtBeginDocument{%
  \providecommand{\repere}{\lycee@repere}%
  \providecommand{\gradx}{\lycee@gradx}%
  \providecommand{\grady}{\lycee@grady}%
  \providecommand{\Cf}{\mathcal{C}\sb{\scriptscriptstyle f}}%
  \providecommand{\Cg}{\mathcal{C}\sb{\scriptscriptstyle g}}%
  \providecommand{\intff}{\lycee@intff}%
  \providecommand{\intfo}{\lycee@intfo}%
  \providecommand{\intof}{\lycee@intof}%
  \providecommand{\intoo}{\lycee@intoo}%
  \providecommand{\ens}{\lycee@ens}}

\newlist{questions}{enumerate}{3}
\setlist[questions,1]{label=\arabic*\textdegree),leftmargin=*,itemsep=2pt,topsep=3pt}
\setlist[questions,2]{label=\alph*),leftmargin=*,itemsep=1pt,topsep=2pt}
\setlist[questions,3]{label=\roman*),leftmargin=*,itemsep=1pt,topsep=1pt}
\setlist[itemize]{label=\textbullet,leftmargin=*,itemsep=2pt,topsep=3pt}

\newcommand{\besoinplace}[1]{\par\penalty\@M\begingroup
  \dimen@=#1\relax\dimen@ii=\pagegoal\advance\dimen@ii-\pagetotal
  \ifdim\dimen@>\dimen@ii\ifdim\dimen@ii>\z@\vfil\fi\break\fi\endgroup}

\newcommand{\lycee@pied}{\thepage/\pageref{LastPage}}
\newcommand{\entete}[2]{%
  \pagestyle{fancy}\fancyhf{}%
  \fancyhead[L]{\footnotesize\color{coulPart}\textbf{#1}}%
  \fancyhead[R]{\footnotesize\color{coulPart}\textbf{#2}}%
  \fancyfoot[C]{\footnotesize\lycee@pied}%
  \renewcommand{\headrulewidth}{0.4pt}%
  \renewcommand{\headrule}{\hbox to\headwidth{%
    \color{coulPart!40}\leaders\hrule height \headrulewidth\hfill}}}

\RequirePackage{ccicons}
\newcommand{\lycee@licence}{{\scriptsize\color{gray}%
  \copyright\ Luc Destombe --- \ccbyncsa\ CC BY-NC-SA 4.0}}
\newif\iflycee@licence \lycee@licencetrue
\newcommand{\sanslicence}{\lycee@licencefalse}
\AtBeginDocument{\iflycee@licence\AddToHookNext{shipout/foreground}{%
  \put(0,0){\rlap{\kern\dimexpr1in+\hoffset+\oddsidemargin+\textwidth\relax
    \llap{\raisebox{-\dimexpr1in+\voffset+\topmargin+\headheight+\headsep
      +\textheight+\footskip\relax}{\lycee@licence}}}}}\fi}

\newcommand{\formulesserrees}{%
  \setlength{\abovedisplayskip}{3pt}\setlength{\belowdisplayskip}{3pt}%
  \setlength{\abovedisplayshortskip}{2pt}\setlength{\belowdisplayshortskip}{3pt}}

\setlength{\parindent}{0pt}

\RequirePackage[implicit=false,hidelinks,pdfencoding=auto,psdextra,bookmarksopen,
  bookmarksopenlevel=1,pdfcreator={},pdfproducer={}]{hyperref}
\RequirePackage{bookmark}
\hypersetup{pdfauthor={Luc Destombe},pdfkeywords={CC BY-NC-SA 4.0}}
\newcounter{lycee@signet}
\newcommand{\signet}[3][]{\stepcounter{lycee@signet}%
  \edef\lycee@dest{\if\relax\detokenize{#1}\relax
    signet.\arabic{lycee@signet}\else#1\fi}%
  \hypertarget{\lycee@dest}{}\bookmark[level=#2,dest=\lycee@dest]{#3}}
\pdfstringdefDisableCommands{%
  \def\R{R}\def\vect#1{#1}\def\Cf{Cf}\def\Cg{Cg}\def\fois#1{×#1}%
  \def\barre#1{#1}\def\intervalle#1#2{[#1;#2]}\def\intff#1#2{[#1;#2]}%
  \def\textsuperscript#1{#1}\def\dfrac#1#2{#1/#2}\def\frac#1#2{#1/#2}%
  \def\vec#1{#1⃗}}%
\RequirePackage[official]{eurosym}

\tcbset{exercice corrige/.style={
  enhanced,breakable,before upper=\raggedright,
  colback=white,colframe=coulExo,
  boxrule=0.8pt,arc=2pt,
  left=6pt,right=6pt,top=8pt,bottom=5pt,
  before skip=9pt,after skip=4pt,
  coltitle=white,fonttitle=\bfseries\small,
  attach boxed title to top left={xshift=8pt,yshift=-\tcboxedtitleheight/2},
  boxed title style={colback=coulExo,arc=2pt,boxrule=0pt}}}
\newcounter{exo}
\newtcolorbox{exercice}[1][]{exercice corrige,
  phantom={\signet[exercice-\arabic{exo}]{2}{Exercice \arabic{exo}}},
  title={Exercice \theexo\ifx\relax#1\relax\else{} --- #1\fi}}
\newenvironment{exo}[1][\relax]{\refstepcounter{exo}\begin{exercice}[#1]}{\end{exercice}}

\geometry{top=1.8cm,bottom=1cm,left=1cm,right=1cm,headheight=14pt,footskip=0.6cm}
\renewcommand{\lycee@pied}{\thepage}

\renewcommand{\theexo}{\ifnum\value{exo}<10 0\fi\arabic{exo}}
\tcbset{exercice corrige/.style={
  enhanced,breakable,before upper=\raggedright,
  colback=white,colframe=coulExo,
  boxrule=0.9pt,arc=2pt,
  left=8pt,right=8pt,top=8pt,bottom=6pt,
  before skip=8pt,after skip=6pt,
  coltitle=white,fonttitle=\bfseries,
  attach boxed title to top left={xshift=8pt,yshift=-\tcboxedtitleheight/2},
  boxed title style={colback=coulExo,arc=2pt,boxrule=0pt}}}

\newcommand{\entetefiche}[2]{%
  \begin{center}
    \begin{tcolorbox}[enhanced,width=0.92\linewidth,colback=white,colframe=coulPart,
        boxrule=1.4pt,arc=3pt,halign=center,top=6pt,bottom=6pt]
      {\large\bfseries\color{coulPart} CORRIGÉ --- FICHE D'EXERCICES}\\[3pt]
      {\normalsize\bfseries #1}\\[2pt]
      {\normalsize #2}
    \end{tcolorbox}
  \end{center}}

\newcounter{partie}
\newcommand{\partie}[1]{%
  \stepcounter{partie}%
  \besoinplace{8\baselineskip}\vspace{4pt}%
  \begin{tcolorbox}[enhanced,colback=coulPart!8,colframe=coulPart,boxrule=0pt,
      leftrule=3pt,arc=0pt,left=6pt,right=4pt,top=3pt,bottom=3pt,
      before skip=6pt,after skip=2pt]
    \bfseries\color{coulPart}\Roman{partie} --- #1\signet{1}{\Roman{partie} --- #1}
  \end{tcolorbox}}
\newcommand{\partiesansnum}[1]{%
  \besoinplace{8\baselineskip}\vspace{4pt}%
  \begin{tcolorbox}[enhanced,colback=coulPart!8,colframe=coulPart,boxrule=0pt,
      leftrule=3pt,arc=0pt,left=6pt,right=4pt,top=3pt,bottom=3pt,
      before skip=6pt,after skip=2pt]
    \bfseries\color{coulPart}#1\signet{1}{#1}
  \end{tcolorbox}}

\setlist[questions,1]{label=\arabic*\textdegree),leftmargin=*,itemsep=3pt,topsep=3pt}
\setlist[questions,2]{label=\alph*),leftmargin=*,itemsep=2pt,topsep=2pt}
\setlist[questions,3]{label=\roman*),leftmargin=*,itemsep=1pt,topsep=1pt}
\newlist{expr}{enumerate}{1}
\setlist[expr]{label={\Alph*\ $=$},leftmargin=*,itemsep=2pt,topsep=3pt}

\newcommand{\res}[1]{{\color{coulRep}#1}}
\newcommand{\rem}[1]{\par\smallskip{\small\itshape\color{black!60}$\rightarrow$ #1}}
\newenvironment{calculs}{\par\smallskip\noindent$\displaystyle\begin{aligned}[t]}%
  {\end{aligned}$\par\smallskip}
\newcommand{\justif}[1]{\quad\text{\small\itshape\color{black!60}#1}}
\newcommand{\cas}[1]{\par\smallskip\noindent\textbf{\color{coulExo}#1}\nobreak}
\newcommand{\explique}[2]{\par\smallskip\noindent
  \begin{minipage}[t]{0.6\linewidth}\raggedright #1\end{minipage}\hfill
  \begin{minipage}[t]{0.37\linewidth}\raggedright\leavevmode
    {\small\itshape\color{black!60}#2}\end{minipage}\par}

\setlength{\parskip}{2pt}

\makeatother
\entete{Chapitre 4 --- Suites numériques --- corrigé}{1\textsuperscript{re} STI2D}
\setlength{\parindent}{0pt}
\setlength{\jot}{5pt}

\let\casligne\cas
\renewcommand{\cas}[1]{\casligne{#1}\par\nopagebreak}

\tcbuselibrary{listings}
\newtcblisting{python}{%
  listing only,enhanced,
  colback=black!3,colframe=black!35,boxrule=0.5pt,arc=1pt,
  left=4pt,right=4pt,top=1pt,bottom=1pt,before skip=3pt,after skip=3pt,
  listing options={language=Python,basicstyle=\ttfamily\small,
    keywordstyle=\color{coulMeth}\bfseries,columns=fullflexible,
    keepspaces=true,showstringspaces=false,aboveskip=0pt,belowskip=0pt}}
\newcommand{\reperesuite}[4]{%
  \draw[grille,step=1] (-1,#2) grid (#1,#3);
  \draw[-{Stealth[length=2mm]},thick] (-1,0)--({#1+0.4},0)
    node[below left=-1pt,font=\footnotesize]{$n$};
  \draw[-{Stealth[length=2mm]},thick] (0,#2)--(0,{#3+0.4})
    node[below left=-1pt,font=\footnotesize]{$#4$};
  \node[below left,font=\scriptsize,inner sep=1.5pt] at (0,0) {$0$};}
\newcommand{\gris}[1]{{\small\itshape\color{black!60}#1}}

\begin{document}

\entetefiche{Chapitre 4 --- Suites numériques}{Première STI2D}

\partie{Définition d'une suite}

\begin{exo}[Suites logiques]
\explique{Pour chaque liste, on calcule l'écart entre deux nombres qui se suivent ; s'il
change, on essaie de multiplier.}{on vérifie la règle sur toute la liste}
\cas{a)}
\explique{$10-5=5$ ; \ $15-10=5$ ; \ $20-15=5$ ; \ $25-20=5$ : on ajoute $5$.}{écart
toujours le même}
\explique{$25+5=30$, puis $35$, puis $40$ : \res{$30$ ; $35$ ; $40$}.}{}
\cas{b)}
\explique{$9-2=7$ ; \ $16-9=7$ ; \ $23-16=7$ ; \ $30-23=7$ : on ajoute $7$.}{}
\explique{$30+7=37$, puis $44$, puis $51$ : \res{$37$ ; $44$ ; $51$}.}{}
\cas{c)}
\explique{$26-30=-4$ ; \ $22-26=-4$ ; \ $18-22=-4$ : on enlève $4$.}{l'écart est
négatif : la liste descend}
\explique{$18-4=14$, puis $10$, puis $6$ : \res{$14$ ; $10$ ; $6$}.}{}
\cas{d)}
\explique{$9-3=6$ mais $27-9=18$ : on n'ajoute pas toujours le même nombre.}{on essaie
de multiplier}
\explique{$\dfrac{9}{3}=3$ ; \ $\dfrac{27}{9}=3$ ; \ $\dfrac{81}{27}=3$ : on multiplie
par $3$.}{}
\explique{$81\times 3=243$, puis $729$, puis $2\,187$ : \res{$243$ ; $729$ ;
$2\,187$}.}{$243\times 3=729$ ; $729\times 3=2\,187$}
\cas{e)}
\explique{Les écarts $7$ ; $19$ ; $37$ ; $61$ changent, et $\dfrac{8}{1}=8$ mais
$\dfrac{27}{8}$ n'est pas entier.}{ni une addition, ni une multiplication}
\explique{Ce sont les cubes : $1^{3}=1$ ; \ $2^{3}=8$ ; \ $3^{3}=27$ ; \ $4^{3}=64$ ; \
$5^{3}=125$.}{$3^{3}=3\times 3\times 3$}
\explique{$6^{3}=216$ ; \ $7^{3}=343$ ; \ $8^{3}=512$ : \res{$216$ ; $343$ ;
$512$}.}{$6\times 6\times 6=36\times 6$}
\cas{f)}
\explique{$1+3=4$ ; \ $3+4=7$ ; \ $4+7=11$ ; \ $7+11=18$ : chaque nombre est la somme des
deux précédents.}{on regarde les deux nombres d'avant}
\explique{$11+18=29$ ; \ $18+29=47$ ; \ $29+47=76$ : \res{$29$ ; $47$ ; $76$}.}{}
\end{exo}

\begin{exo}[Rang et position]
\explique{D'un terme au suivant, on ajoute $6$ : $13-7=6$ ; $19-13=6$ ; $25-19=6$. On
prolonge la liste en numérotant à partir de $u_{0}$ :}{le premier terme est $u_{0}$}

\smallskip
{\renewcommand{\arraystretch}{1.3}
\begin{tabular}{|c|c|c|c|c|c|c|c|c|c|}
\hline
Position & $1^{\text{er}}$ & $2^{\text{e}}$ & $3^{\text{e}}$ & $4^{\text{e}}$ & $5^{\text{e}}$
  & $6^{\text{e}}$ & $7^{\text{e}}$ & $8^{\text{e}}$ & $9^{\text{e}}$\\
\hline
Rang $n$ & $0$ & $1$ & $2$ & $3$ & $4$ & $5$ & $6$ & $7$ & $8$\\
\hline
$u_{n}$ & $7$ & $13$ & $19$ & $25$ & $31$ & $37$ & $43$ & $49$ & $55$\\
\hline
\end{tabular}}
\cas{1)}
\explique{$u_{1}=\res{13}$ et $u_{3}=\res{25}$.}{$u_{1}$ est le $2^{\text{e}}$ nombre de
la liste}
\cas{2)}
\explique{Le $5^{\text{e}}$ terme est de rang \res{$4$} : c'est $u_{4}=25+6=\res{31}$.}{le
rang vaut la position moins $1$}
\cas{3)}
\explique{$u_{8}=\res{55}$ : c'est le \res{$9^{\text{e}}$} nombre de la liste.}{la
position vaut le rang plus $1$}
\rem{Erreur fréquente : prendre $u_{4}$ pour le $4^{\text{e}}$ terme. Comme on commence à
$u_{0}$, le $4^{\text{e}}$ terme est $u_{3}$.}
\end{exo}

\begin{exo}[Lire un tableau]
\cas{1)}
\explique{Dans la colonne $n=3$, on lit $u_{3}=\res{3}$.}{le rang est dans la première
ligne}
\explique{Le $3^{\text{e}}$ terme est $u_{2}$ : il vaut \res{$7$}.}{$1^{\text{er}}$ terme :
$u_{0}$ ; $2^{\text{e}}$ : $u_{1}$ ; $3^{\text{e}}$ : $u_{2}$}
\cas{2)}
\explique{On cherche $-1$ dans la deuxième ligne : il est dans la colonne $n=4$. Donc
$u_{n}=-1$ pour \res{$n=4$}.}{on lit le tableau dans l'autre sens}
\cas{3)}
\explique{$11-15=-4$ ; \ $7-11=-4$ ; \ $3-7=-4$ ; \ $-1-3=-4$ ; \ $-5-(-1)=-4$ : on
enlève $4$ à chaque fois.}{$-5-(-1)=-5+1$}
\begin{calculs}
u_{6} &= u_{5}-4 &&\justif{on enlève $4$ au terme précédent}\\
u_{6} &= -5-4 = \res{-9} &&
\end{calculs}
\end{exo}

\begin{exo}[Les tribunes d'un stade]
\cas{1)}
\begin{calculs}
u_{1} &= 30+4 = \res{34} &&\justif{la $2^{\text{e}}$ rangée : $4$ sièges de plus}\\
u_{2} &= 34+4 = \res{38} &&\justif{toujours à partir de la rangée d'avant}\\
u_{3} &= 38+4 = \res{42} &&\\
u_{4} &= 42+4 = \res{46} &&
\end{calculs}
\cas{2)}
\explique{La $10^{\text{e}}$ rangée correspond à $u_{9}$.}{la $1^{\text{re}}$ rangée est
$u_{0}$ : le rang vaut la position moins $1$}
\begin{calculs}
u_{5} &= 46+4 = 50 &&\justif{on continue la liste}\\
u_{6} &= 50+4 = 54 &&\\
u_{7} &= 54+4 = 58 &&\\
u_{8} &= 58+4 = 62 &&\\
u_{9} &= 62+4 = 66 &&
\end{calculs}
\explique{La $10^{\text{e}}$ rangée compte \res{$66$ sièges}.}{contrôle : de $u_{0}$ à
$u_{9}$, on a ajouté $9$ fois $4$ ; $30+36=66$}
\cas{3)}
\explique{De $30$ à $70$, il faut $40$ sièges de plus, soit $10$ fois $4$ : c'est
$u_{10}$.}{$40\div 4=10$}
\explique{Contrôle : $u_{10}=66+4=70$.}{on ajoute $4$ à $u_{9}$}
\explique{$u_{10}$ est la \res{$11^{\text{e}}$ rangée}.}{la position vaut le rang plus
$1$}
\rem{Piège du décalage : $u_{10}$ n'est pas la $10^{\text{e}}$ rangée, mais la
$11^{\text{e}}$.}
\end{exo}

\partie{Suites définies explicitement}

\begin{exo}[Des entiers]
\begin{calculs}
u_{0} &= 5\times 0-8 &&\justif{on remplace $n$ par $0$}\\
u_{0} &= \res{-8} &&\\
u_{1} &= 5\times 1-8 = \res{-3} &&\justif{on remplace $n$ par $1$}\\
u_{2} &= 5\times 2-8 = \res{2} &&\justif{$10-8$}\\
u_{3} &= 5\times 3-8 = \res{7} &&\justif{$15-8$}\\
u_{4} &= 5\times 4-8 = \res{12} &&\justif{$20-8$ ; d'un terme au suivant, on ajoute $5$}\\
u_{20} &= 5\times 20-8 &&\justif{directement, sans les termes d'avant}\\
u_{20} &= 100-8 = \res{92} &&\\
u_{100} &= 5\times 100-8 &&\\
u_{100} &= 500-8 = \res{492} &&
\end{calculs}
\end{exo}

\begin{exo}[Pour chaque formule]
\explique{Les quatre premiers termes sont ceux de rang $0$ ; $1$ ; $2$ et $3$.}{le premier
terme est de rang $0$}
\cas{a)}
\begin{calculs}
a_{0} &= 3^{0} = \res{1} &&\justif{un nombre à la puissance $0$ vaut $1$}\\
a_{1} &= 3^{1} = \res{3} &&\\
a_{2} &= 3^{2} = \res{9} &&\justif{$3\times 3$, et non $3\times 2$}\\
a_{3} &= 3^{3} = \res{27} &&\justif{$9\times 3$}\\
a_{10} &= 3^{10} = \res{59\,049} &&\justif{à la calculatrice}
\end{calculs}
\cas{b)}
\begin{calculs}
b_{0} &= 0^{2}+0 = \res{0} &&\justif{on remplace les deux $n$ par $0$}\\
b_{1} &= 1^{2}+1 = \res{2} &&\\
b_{2} &= 2^{2}+2 = \res{6} &&\justif{$4+2$}\\
b_{3} &= 3^{2}+3 = \res{12} &&\justif{$9+3$}\\
b_{10} &= 10^{2}+10 = \res{110} &&\justif{$100+10$}
\end{calculs}
\cas{c)}
\begin{calculs}
c_{0} &= -3\times 0+4 = \res{4} &&\\
c_{1} &= -3\times 1+4 = \res{1} &&\justif{$-3+4$}\\
c_{2} &= -3\times 2+4 = \res{-2} &&\justif{$-6+4$}\\
c_{3} &= -3\times 3+4 = \res{-5} &&\justif{$-9+4$}\\
c_{10} &= -3\times 10+4 = \res{-26} &&\justif{$-30+4$}
\end{calculs}
\cas{d)}
\begin{calculs}
d_{0} &= 0^{2}-4\times 0 = \res{0} &&\\
d_{1} &= 1^{2}-4\times 1 = \res{-3} &&\justif{$1-4$}\\
d_{2} &= 2^{2}-4\times 2 = \res{-4} &&\justif{$4-8$}\\
d_{3} &= 3^{2}-4\times 3 = \res{-3} &&\justif{$9-12$}\\
d_{10} &= 10^{2}-4\times 10 = \res{60} &&\justif{$100-40$}
\end{calculs}
\cas{e)}
\begin{calculs}
e_{0} &= \dfrac{1}{2}\times 0^{2}-1 = \res{-1} &&\\
e_{1} &= \dfrac{1}{2}\times 1^{2}-1 = \res{-\dfrac{1}{2}} &&\justif{$\dfrac{1}{2}-1$}\\
e_{2} &= \dfrac{1}{2}\times 2^{2}-1 = \res{1} &&\justif{le carré d'abord : $\dfrac{1}{2}\times 4=2$}\\
e_{3} &= \dfrac{1}{2}\times 3^{2}-1 &&\justif{$3^{2}=9$}\\
e_{3} &= \dfrac{9}{2}-\dfrac{2}{2} = \res{\dfrac{7}{2}} &&\justif{$1=\dfrac{2}{2}$}\\
e_{10} &= \dfrac{1}{2}\times 10^{2}-1 &&\\
e_{10} &= 50-1 = \res{49} &&\justif{la moitié de $100$}
\end{calculs}
\rem{Erreur fréquente : calculer $\left(\dfrac{1}{2}\times 3\right)^{2}$. Le carré ne porte
que sur $n$ : on calcule $3^{2}$, puis on prend la moitié.}
\cas{f)}
\begin{calculs}
f_{0} &= (0-3)^{2} = (-3)^{2} &&\justif{la parenthèse d'abord}\\
f_{0} &= \res{9} &&\justif{$(-3)\times(-3)$}\\
f_{1} &= (1-3)^{2} = (-2)^{2} &&\\
f_{1} &= \res{4} &&\\
f_{2} &= (2-3)^{2} = \res{1} &&\justif{$(-1)^{2}=1$}\\
f_{3} &= (3-3)^{2} = \res{0} &&\\
f_{10} &= (10-3)^{2} = 7^{2} &&\\
f_{10} &= \res{49} &&
\end{calculs}
\end{exo}

\begin{exo}[Fractions et signes]
\cas{1)}
\begin{calculs}
u_{0} &= \dfrac{0+1}{0+2} = \res{\dfrac{1}{2}} &&\justif{on remplace $n$ par $0$, en
haut et en bas}\\
u_{1} &= \dfrac{1+1}{1+2} = \res{\dfrac{2}{3}} &&\\
u_{2} &= \dfrac{2+1}{2+2} = \res{\dfrac{3}{4}} &&\\
u_{3} &= \dfrac{3+1}{3+2} = \res{\dfrac{4}{5}} &&\\
u_{8} &= \dfrac{8+1}{8+2} = \res{\dfrac{9}{10}} &&\justif{aucune ne se simplifie}
\end{calculs}
\cas{2)}
\begin{calculs}
v_{1} &= (-1)^{1} = \res{-1} &&\\
v_{2} &= (-1)^{2} = \res{1} &&\justif{$(-1)\times(-1)=1$}\\
v_{3} &= (-1)^{3} = \res{-1} &&\justif{$1\times(-1)$}\\
v_{4} &= (-1)^{4} = \res{1} &&
\end{calculs}
\explique{Chaque facteur $-1$ change le signe : au rang pair, $v_{n}=1$ ; au rang impair,
$v_{n}=-1$.}{les $-1$ se regroupent par deux}
\explique{$40$ est pair : $v_{40}=\res{1}$. \ $65$ est impair : $v_{65}=\res{-1}$. \
$100$ est pair : $v_{100}=\res{1}$.}{un nombre pair se termine par $0$ ; $2$ ; $4$ ;
$6$ ou $8$}
\end{exo}

\begin{exo}[$u_{n+1}$ ou $u_{n}+1$ ?]
\explique{$u_{n+1}$ est le terme suivant : on remplace $n$ par $n+1$, entre parenthèses.
$u_{n}+1$ est le terme $u_{n}$ augmenté de $1$.}{ce ne sont pas les mêmes nombres}
\cas{a)}
\begin{calculs}
u_{n+1} &= 3(n+1)+2 &&\justif{on remplace $n$ par $n+1$}\\
u_{n+1} &= 3n+3+2 &&\justif{on développe}\\
u_{n+1} &= \res{3n+5} &&\\
u_{n}+1 &= 3n+2+1 &&\justif{on ajoute $1$ à $u_{n}$}\\
u_{n}+1 &= \res{3n+3} &&
\end{calculs}
\explique{Contrôle avec $n=1$ : $u_{2}=3\times 2+2=8$, et $3\times 1+5=8$.}{on vérifie la
formule de $u_{n+1}$ sur un exemple}
\cas{b)}
\begin{calculs}
u_{n+1} &= 5-4(n+1) &&\\
u_{n+1} &= 5-4n-4 &&\justif{$-4\times(+1)=-4$}\\
u_{n+1} &= \res{1-4n} &&\\
u_{n}+1 &= 5-4n+1 &&\\
u_{n}+1 &= \res{6-4n} &&
\end{calculs}
\cas{c)}
\begin{calculs}
u_{n+1} &= (n+1)^{2}-2 &&\\
u_{n+1} &= n^{2}+2n+1-2 &&\justif{identité remarquable $(a+b)^{2}$}\\
u_{n+1} &= \res{n^{2}+2n-1} &&\\
u_{n}+1 &= n^{2}-2+1 &&\\
u_{n}+1 &= \res{n^{2}-1} &&
\end{calculs}
\rem{Erreurs fréquentes : $3(n+1)=3n+1$ (le $3$ multiplie aussi le $1$) ;
$(n+1)^{2}=n^{2}+1$ (il manque $2n$).}
\end{exo}

\begin{exo}[La borne de recharge]
\cas{1)}
\explique{$p_{n}=\res{2+3n}$}{la connexion, plus $3$~€ pour chacune des $n$ heures}
\cas{2)}
\begin{calculs}
p_{5} &= 2+3\times 5 &&\justif{on remplace $n$ par $5$}\\
p_{5} &= 2+15 = \res{17} &&\\
p_{8} &= 2+3\times 8 &&\\
p_{8} &= 2+24 = \res{26} &&
\end{calculs}
\explique{$5$ heures de charge coûtent $17$~€ ; $8$ heures coûtent $26$~€.}{}
\cas{3)}
\explique{On cherche $n$ tel que $p_{n}=23$ : on résout une équation dont l'inconnue est
le nombre d'heures.}{on connaît le prix, on cherche $n$}
\begin{calculs}
2+3n &= 23 &&\justif{on traduit $p_{n}=23$}\\
3n &= 21 &&\justif{on enlève $2$ des deux côtés}\\
n &= 7 &&\justif{on divise par $3$}
\end{calculs}
\explique{Sa voiture a chargé \res{$7$ heures}.}{contrôle : $p_{7}=2+21=23$}
\end{exo}

\partie{Suites définies par récurrence}

\begin{exo}[Tripler et enlever]
\begin{calculs}
u_{1} &= 3u_{0}-2 &&\justif{on remplace $n$ par $0$ dans la relation}\\
u_{1} &= 3\times 2-2 = \res{4} &&\justif{on remplace $u_{0}$ par sa valeur}\\
u_{2} &= 3\times 4-2 = \res{10} &&\justif{$u_{2}=3u_{1}-2$}\\
u_{3} &= 3\times 10-2 = \res{28} &&\justif{toujours le terme d'avant}\\
u_{4} &= 3\times 28-2 = \res{82} &&\justif{$3\times 28=84$}\\
u_{5} &= 3\times 82-2 = \res{244} &&\justif{$3\times 82=246$}
\end{calculs}
\rem{On ne peut pas calculer $u_{5}$ directement : il faut $u_{4}$, donc $u_{3}$\ldots{}
Les termes se calculent dans l'ordre.}
\end{exo}

\begin{exo}[Des relatifs]
\cas{a)}
\begin{calculs}
u_{1} &= -2u_{0}+3 &&\justif{la relation avec $n=0$}\\
u_{1} &= -2\times 4+3 = \res{-5} &&\justif{$-8+3$}\\
u_{2} &= -2\times(-5)+3 &&\justif{on remplace $u_{1}$ par $-5$, entre parenthèses}\\
u_{2} &= 10+3 = \res{13} &&\justif{moins par moins donne plus}\\
u_{3} &= -2\times 13+3 = \res{-23} &&\justif{$-26+3$}\\
u_{4} &= -2\times(-23)+3 &&\\
u_{4} &= 46+3 = \res{49} &&
\end{calculs}
\cas{b)}
\begin{calculs}
v_{1} &= v_{0}^{2}-3 &&\justif{la relation avec $n=0$}\\
v_{1} &= 1^{2}-3 = \res{-2} &&\\
v_{2} &= (-2)^{2}-3 &&\justif{on remplace $v_{1}$ par $-2$, entre parenthèses}\\
v_{2} &= 4-3 = \res{1} &&\justif{$(-2)\times(-2)=4$}\\
v_{3} &= 1^{2}-3 = \res{-2} &&\justif{même calcul que pour $v_{1}$}\\
v_{4} &= (-2)^{2}-3 = \res{1} &&
\end{calculs}
\explique{Les termes alternent : $1$ ; $-2$ ; $1$ ; $-2$\ldots}{dès qu'on retombe sur $1$,
tout recommence}
\cas{c)}
\begin{calculs}
w_{1} &= 6-w_{0} &&\justif{la relation avec $n=0$}\\
w_{1} &= 6-(-4) &&\justif{on remplace $w_{0}$ par $-4$, entre parenthèses}\\
w_{1} &= 6+4 = \res{10} &&\justif{enlever $-4$, c'est ajouter $4$}\\
w_{2} &= 6-10 = \res{-4} &&\\
w_{3} &= 6-(-4) = \res{10} &&\justif{même calcul que pour $w_{1}$}\\
w_{4} &= 6-10 = \res{-4} &&
\end{calculs}
\rem{Erreur fréquente : écrire $-2\times-5$ ou $-2^{2}$. Un nombre négatif se remplace
toujours entre parenthèses : $(-2)^{2}=4$, alors que $-2^{2}=-4$.}
\end{exo}

\begin{exo}[Sans tout calculer ?]
\cas{1)}
\begin{calculs}
a_{1} &= -a_{0} &&\justif{la relation avec $n=0$}\\
a_{1} &= -(-3) = \res{3} &&\justif{on prend l'opposé}\\
a_{2} &= -a_{1} = \res{-3} &&\justif{l'opposé de $3$}\\
a_{3} &= -(-3) = \res{3} &&\\
a_{4} &= \res{-3} &&\\
a_{5} &= \res{3} &&
\end{calculs}
\explique{Au rang pair, on trouve $-3$ (comme $a_{0}$) ; au rang impair, $3$.}{chaque
étape change le signe}
\explique{$40$ est pair : $a_{40}=\res{-3}$. \ $77$ est impair : $a_{77}=\res{3}$.}{}
\cas{2)}
\begin{calculs}
b_{1} &= b_{0}+4 &&\justif{la relation avec $n=0$}\\
b_{1} &= 5+4 = \res{9} &&\\
b_{2} &= 9+4 = \res{13} &&\\
b_{3} &= 13+4 = \res{17} &&\\
b_{4} &= 17+4 = \res{21} &&\\
b_{5} &= 21+4 = \res{25} &&
\end{calculs}
\explique{Chaque terme demande un calcul : de $b_{1}$ à $b_{50}$, il faut \res{$50$
calculs}.}{un calcul par terme}
\rem{Une formule explicite donnerait $b_{50}$ en un seul calcul : on la verra avec les
suites arithmétiques.}
\end{exo}

\begin{exo}[Avec des fractions]
\begin{calculs}
u_{1} &= \dfrac{10}{u_{0}} &&\justif{la relation avec $n=0$}\\
u_{1} &= \dfrac{10}{50} = \dfrac{1\times\barre{10}}{5\times\barre{10}} &&\justif{on
simplifie par $10$}\\
u_{1} &= \res{\dfrac{1}{5}} &&\\
u_{2} &= 10\div\dfrac{1}{5} &&\justif{$u_{2}=\dfrac{10}{u_{1}}$}\\
u_{2} &= 10\times 5 = \res{50} &&\justif{diviser par $\dfrac{1}{5}$, c'est multiplier par
$5$}\\
u_{3} &= \dfrac{10}{50} = \res{\dfrac{1}{5}} &&\justif{même calcul que pour $u_{1}$}\\
u_{4} &= \res{50} &&\justif{même calcul que pour $u_{2}$}
\end{calculs}
\explique{On remarque que les termes \res{alternent entre $50$ et $\dfrac{1}{5}$}.}{on
retombe sur $u_{0}$ : tout recommence}
\end{exo}

\begin{exo}[Quand la relation contient $n$]
\explique{Pour calculer $u_{n+1}$, on remplace $n$ par le rang du terme
\textbf{précédent}.}{pour $u_{1}$ : $n=0$ ; pour $u_{2}$ : $n=1$\ldots}
\begin{calculs}
u_{1} &= u_{0}-2\times 0+5 &&\justif{$n=0$}\\
u_{1} &= 2-0+5 = \res{7} &&\\
u_{2} &= u_{1}-2\times 1+5 &&\justif{$n=1$}\\
u_{2} &= 7-2+5 = \res{10} &&\\
u_{3} &= u_{2}-2\times 2+5 &&\justif{$n=2$}\\
u_{3} &= 10-4+5 = \res{11} &&\\
u_{4} &= u_{3}-2\times 3+5 &&\justif{$n=3$}\\
u_{4} &= 11-6+5 = \res{10} &&
\end{calculs}
\rem{Erreur fréquente : prendre $n=1$ pour calculer $u_{1}$. Dans $u_{n+1}$, si le rang
est $1$, alors $n=0$.}
\end{exo}

\begin{exo}[Une culture de levures]
\cas{1)}
\explique{$b_{0}=\res{300}$}{le nombre de levures au départ, après $0$ heure}
\explique{$b_{n+1}=\res{2b_{n}-80}$}{d'une heure à la suivante : on double (on multiplie par
$2$), puis on enlève $80$}
\cas{2)}
\begin{calculs}
b_{1} &= 2\times 300-80 &&\justif{$b_{1}=2b_{0}-80$}\\
b_{1} &= 600-80 = \res{520} &&\\
b_{2} &= 2\times 520-80 &&\\
b_{2} &= 1\,040-80 = \res{960} &&\\
b_{3} &= 2\times 960-80 &&\\
b_{3} &= 1\,920-80 = \res{1\,840} &&
\end{calculs}
\explique{Après $3$ heures, la cuve contient $1\,840$~milliers de levures, soit
$1\,840\,000$ levures.}{les termes sont en milliers}
\rem{L'ordre compte : on double \emph{puis} on prélève. Écrire $2(b_{n}-80)$ reviendrait
à prélever avant de doubler.}
\end{exo}

\partie{Programmer une suite définie par récurrence}

\begin{exo}[Que renvoie ce programme ?]
\cas{1)}
\explique{La ligne \texttt{u = 3} donne le premier terme : $u_{0}=\res{3}$.}{avant la
boucle : le premier terme}
\explique{La ligne de la boucle \texttt{u = 3*u - 4} donne la relation :
$u_{n+1}=\res{3u_{n}-4}$.}{le nouveau \texttt{u} est $u_{n+1}$, l'ancien est $u_{n}$}
\cas{2)}
\explique{\texttt{suite(n)} fait tourner la boucle $n$ fois à partir de $u_{0}$ : elle
renvoie $u_{n}$.}{\texttt{suite(5)} applique $5$ fois la relation}
\begin{calculs}
u_{1} &= 3\times 3-4 = 5 &&\justif{$u_{1}=3u_{0}-4$}\\
u_{2} &= 3\times 5-4 = 11 &&\\
u_{3} &= 3\times 11-4 = 29 &&\justif{$33-4$}\\
u_{4} &= 3\times 29-4 = 83 &&\justif{$87-4$}\\
u_{5} &= 3\times 83-4 = 245 &&\justif{$249-4$}
\end{calculs}
\explique{\texttt{suite(1)} renvoie \res{$5$}, \texttt{suite(2)} renvoie \res{$11$} et
\texttt{suite(5)} renvoie \res{$245$}.}{}
\rem{Erreur fréquente : croire que \texttt{suite(1)} renvoie $u_{0}$. La boucle tourne une
fois : elle renvoie $u_{1}$.}
\end{exo}

\begin{exo}[Même travail]
\cas{1)}
\explique{$u_{0}=\res{160}$ et $u_{n+1}=\res{0{,}5\,u_{n}+2}$.}{\texttt{u = 160} avant la
boucle ; \texttt{u = 0.5*u + 2} dans la boucle}
\cas{2)}
\begin{calculs}
u_{1} &= 0{,}5\times 160+2 &&\justif{$0{,}5\times 160$ : la moitié de $160$}\\
u_{1} &= 80+2 = 82 &&\\
u_{2} &= 0{,}5\times 82+2 &&\\
u_{2} &= 41+2 = 43 &&\\
u_{3} &= 0{,}5\times 43+2 &&\\
u_{3} &= 21{,}5+2 = 23{,}5 &&\\
u_{4} &= 0{,}5\times 23{,}5+2 &&\\
u_{4} &= 11{,}75+2 = 13{,}75 &&\\
u_{5} &= 0{,}5\times 13{,}75+2 &&\\
u_{5} &= 6{,}875+2 = 8{,}875 &&\\
u_{6} &= 0{,}5\times 8{,}875+2 &&\\
u_{6} &= 4{,}4375+2 = 6{,}4375 &&
\end{calculs}
\explique{\texttt{suite(1)} renvoie \res{$82$}, \texttt{suite(2)} renvoie \res{$43$} et
\texttt{suite(6)} renvoie \res{$6{,}4375$}.}{Python affiche \texttt{82.0} : avec
\texttt{0.5}, le résultat est un nombre décimal}
\end{exo}

\begin{exo}[Compléter le programme]
\cas{1)}
\begin{calculs}
u_{1} &= 1+2u_{0} &&\justif{la relation avec $n=0$}\\
u_{1} &= 1+2\times 2 = 5 &&\\
u_{2} &= 1+2\times 5 = 11 &&\\
u_{3} &= 1+2\times 11 = \res{23} &&
\end{calculs}
\cas{2)}
\begin{minipage}[t]{0.42\linewidth}\vspace{0pt}
\begin{python}
def suite(n):
    u = 2
    for i in range(n):
        u = 1 + 2*u
    return u
\end{python}
\end{minipage}\hfill
\begin{minipage}[t]{0.54\linewidth}\vspace{0pt}\raggedright
\gris{la ligne \texttt{u = 2} donne $u_{0}$ ; dans la boucle, on traduit
$u_{n+1}=1+2u_{n}$ : le nouveau \texttt{u} se calcule avec l'ancien}
\end{minipage}
\explique{Dans la console, \texttt{suite(10)} affiche \texttt{3071} :
$u_{10}=\res{3\,071}$.}{contrôle : \texttt{suite(3)} affiche \texttt{23}, comme à la
question 1)}
\cas{3)}
\begin{minipage}[t]{0.42\linewidth}\vspace{0pt}
\begin{python}
def suite(n):
    u = 1
    for i in range(n):
        u = 6 - 2*u
    return u
\end{python}
\end{minipage}\hfill
\begin{minipage}[t]{0.54\linewidth}\vspace{0pt}\raggedright
\gris{seules la valeur de départ ($v_{0}=1$) et la ligne de la boucle changent}
\end{minipage}
\begin{calculs}
v_{1} &= 6-2\times 1 = \res{4} &&\justif{à la main, pour contrôler le programme}
\end{calculs}
\explique{\texttt{suite(5)} affiche \texttt{34} et \texttt{suite(10)} affiche
\texttt{-1022} : $v_{5}=\res{34}$ et $v_{10}=\res{-1\,022}$.}{\texttt{suite(1)} doit
afficher \texttt{4}}
\rem{Erreur fréquente : écrire \texttt{u = 6 - 2*v}. Dans le programme, la suite s'appelle
toujours \texttt{u} : il n'y a pas de variable \texttt{v}.}
\end{exo}

\begin{exo}[Les abonnés d'une salle de sport]
\cas{1)}
\explique{Si $40\,\%$ des abonnés partent, $60\,\%$ restent : on multiplie par $0{,}6$,
puis on ajoute les $600$ nouveaux.}{$100\,\%-40\,\%=60\,\%$, et $60\,\%=0{,}6$}
\begin{calculs}
a_{1} &= 0{,}6\times 1\,200+600 &&\justif{les abonnés restés, plus les nouveaux}\\
a_{1} &= 720+600 = \res{1\,320} &&\\
a_{2} &= 0{,}6\times 1\,320+600 &&\\
a_{2} &= 792+600 = \res{1\,392} &&
\end{calculs}
\cas{2)}
\explique{$a_{n+1}=\res{0{,}6\,a_{n}+600}$}{le même calcul, d'une année à la suivante}
\cas{3)}
\begin{minipage}[t]{0.42\linewidth}\vspace{0pt}
\begin{python}
def abonnes(n):
    a = 1200
    for i in range(n):
        a = 0.6*a + 600
    return a
\end{python}
\end{minipage}\hfill
\begin{minipage}[t]{0.54\linewidth}\vspace{0pt}\raggedright
\gris{sur le modèle de \texttt{suite(n)} : \texttt{a = 1200} donne $a_{0}$ ; la boucle
traduit la relation de la question 2)}
\end{minipage}
\explique{Contrôle : \texttt{abonnes(2)} affiche \texttt{1392.0}, comme à la question
1).}{}
\cas{4)}
\explique{2039 est l'année $2024+15$ : on cherche $a_{15}$.}{on compte $15$ ans après 2024}
\explique{\texttt{abonnes(15)} affiche \texttt{1499.85\ldots} : en 2039, il y aura
environ \res{$1\,500$ abonnés}.}{arrondi à l'unité}
\cas{5)}
\explique{Le nombre d'abonnés semble se stabiliser vers \res{$1\,500$}.}{on essaie
\texttt{abonnes(20)}, \texttt{abonnes(30)} : on trouve presque $1\,500$}
\rem{Non exigé : avec $1\,500$ abonnés, $0{,}6\times 1\,500+600=900+600=1\,500$ ; le
nombre ne change plus.}
\end{exo}

\partie{Représentation graphique d'une suite}

\begin{exo}[Placer les points]
\cas{1)}
\begin{minipage}[t]{0.58\linewidth}\raggedright\vspace{0pt}
\begin{calculs}
u_{0} &= 0^{2}-6\times 0+5 = 5 &&\justif{on remplace $n$ par $0$}\\
u_{1} &= 1^{2}-6\times 1+5 = 0 &&\justif{$1-6+5$}\\
u_{2} &= 2^{2}-6\times 2+5 = -3 &&\justif{$4-12+5$}\\
u_{3} &= 3^{2}-6\times 3+5 = -4 &&\justif{$9-18+5$}\\
u_{4} &= 4^{2}-6\times 4+5 = -3 &&\justif{$16-24+5$}\\
u_{5} &= 5^{2}-6\times 5+5 = 0 &&\justif{$25-30+5$}\\
u_{6} &= 6^{2}-6\times 6+5 = 5 &&\justif{$36-36+5$}
\end{calculs}

{\renewcommand{\arraystretch}{1.3}
\begin{tabular}{|c|c|c|c|c|c|c|c|}
\hline
$n$ & $0$ & $1$ & $2$ & $3$ & $4$ & $5$ & $6$\\
\hline
$u_{n}$ & \res{$5$} & \res{$0$} & \res{$-3$} & \res{$-4$} & \res{$-3$} & \res{$0$}
  & \res{$5$}\\
\hline
\end{tabular}}

\cas{2)}
\explique{Chaque colonne donne un point $(n\,;u_{n})$ : par exemple, une croix en
$(3\,;-4)$.}{le rang en abscisse, le terme en ordonnée}
\explique{Les croix ne sont pas reliées.}{il n'y a pas de rang $2{,}5$}
\end{minipage}\hfill
\begin{minipage}[t]{0.38\linewidth}\vspace{0pt}\centering
\begin{tikzpicture}[x=0.5cm,y=0.5cm]
  \reperesuite{7}{-5}{6}{u_{n}}
  \gradx{1,2,3,4,5,6} \grady{-4,-3,-2,-1,1,2,3,4,5}
  \foreach \p in {(0,5),(1,0),(2,-3),(3,-4),(4,-3),(5,0),(6,5)}
    \path[coulCourbe] \p node[croix]{};
\end{tikzpicture}
\end{minipage}
\end{exo}

\begin{exo}[Une suite qui alterne]
\begin{minipage}[t]{0.58\linewidth}\raggedright\vspace{0pt}
\begin{calculs}
u_{1} &= -u_{0}+2 &&\justif{la relation avec $n=0$}\\
u_{1} &= -4+2 = -2 &&\\
u_{2} &= -(-2)+2 &&\justif{l'opposé de $-2$, entre parenthèses}\\
u_{2} &= 2+2 = 4 &&\justif{on retrouve $u_{0}$}\\
u_{3} &= -4+2 = -2 &&\justif{tout recommence}\\
u_{4} &= 4 &&\\
u_{5} &= -2 &&\\
u_{6} &= 4 &&
\end{calculs}

{\renewcommand{\arraystretch}{1.3}
\begin{tabular}{|c|c|c|c|c|c|c|c|}
\hline
$n$ & $0$ & $1$ & $2$ & $3$ & $4$ & $5$ & $6$\\
\hline
$u_{n}$ & $4$ & \res{$-2$} & \res{$4$} & \res{$-2$} & \res{$4$} & \res{$-2$} & \res{$4$}\\
\hline
\end{tabular}}

\explique{Les croix alternent entre deux niveaux : $4$ aux rangs pairs, $-2$ aux rangs
impairs.}{on place les points $(n\,;u_{n})$}
\end{minipage}\hfill
\begin{minipage}[t]{0.38\linewidth}\vspace{0pt}\centering
\begin{tikzpicture}[x=0.5cm,y=0.5cm]
  \reperesuite{7}{-3}{5}{u_{n}}
  \gradx{1,2,3,4,5,6} \grady{-2,-1,1,2,3,4}
  \foreach \p in {(0,4),(1,-2),(2,4),(3,-2),(4,4),(5,-2),(6,4)}
    \path[coulCourbe] \p node[croix]{};
\end{tikzpicture}
\end{minipage}
\end{exo}

\begin{exo}[Une suite qui se rapproche]
\begin{minipage}[t]{0.58\linewidth}\raggedright\vspace{0pt}
\begin{calculs}
u_{1} &= \dfrac{u_{0}+10}{2} &&\justif{la relation avec $n=0$}\\
u_{1} &= \dfrac{-6+10}{2} = \dfrac{4}{2} &&\justif{on remplace $u_{0}$ par $-6$}\\
u_{1} &= 2 &&\\
u_{2} &= \dfrac{2+10}{2} = \dfrac{12}{2} &&\\
u_{2} &= 6 &&\\
u_{3} &= \dfrac{6+10}{2} = \dfrac{16}{2} &&\\
u_{3} &= 8 &&\\
u_{4} &= \dfrac{8+10}{2} = \dfrac{18}{2} &&\\
u_{4} &= 9 &&
\end{calculs}

{\renewcommand{\arraystretch}{1.3}
\begin{tabular}{|c|c|c|c|c|c|}
\hline
$n$ & $0$ & $1$ & $2$ & $3$ & $4$\\
\hline
$u_{n}$ & $-6$ & \res{$2$} & \res{$6$} & \res{$8$} & \res{$9$}\\
\hline
\end{tabular}}

\explique{Les termes semblent se rapprocher de \res{$10$}.}{les écarts à $10$ : $16$ ;
$8$ ; $4$ ; $2$ ; $1$, divisés par $2$ à chaque étape}
\end{minipage}\hfill
\begin{minipage}[t]{0.38\linewidth}\vspace{0pt}\centering
\begin{tikzpicture}[x=0.4cm,y=0.3cm]
  \reperesuite{5}{-7}{11}{u_{n}}
  \gradx{1,2,3,4} \grady{-6,-4,-2,2,4,6,8,10}
  \draw[coulPart,dashed] (-1,10)--(5,10);
  \foreach \p in {(0,-6),(1,2),(2,6),(3,8),(4,9)}
    \path[coulCourbe] \p node[croix]{};
\end{tikzpicture}
\end{minipage}
\rem{Chaque terme est le milieu entre le terme précédent et $10$ : il se rapproche de $10$
sans le dépasser.}
\end{exo}

\begin{exo}[Lire un graphique]
\begin{minipage}[t]{0.58\linewidth}\raggedright\vspace{0pt}
\cas{1)}
\explique{La croix d'abscisse $0$ est à la hauteur $8$ : $u_{0}=\res{8}$.}{l'ordonnée de
la croix d'abscisse $n$ est $u_{n}$}
\explique{De même : $u_{1}=\res{6}$ ; \ $u_{2}=\res{4}$ ; \ $u_{3}=\res{2}$ ; \
$u_{4}=\res{0}$.}{$u_{4}=0$ : la croix est sur l'axe}
\cas{2)}
\explique{$6-8=-2$ ; \ $4-6=-2$ ; \ $2-4=-2$ ; \ $0-2=-2$ : on \res{enlève $2$} à chaque
fois.}{d'une croix à la suivante, on descend de $2$ carreaux}
\cas{3)}
\begin{calculs}
u_{5} &= u_{4}-2 &&\justif{on enlève $2$}\\
u_{5} &= 0-2 = \res{-2} &&\\
u_{6} &= -2-2 = \res{-4} &&
\end{calculs}
\end{minipage}\hfill
\begin{minipage}[t]{0.38\linewidth}\vspace{0pt}\centering
\begin{tikzpicture}[x=0.5cm,y=0.5cm]
  \reperesuite{7}{-5}{9}{u_{n}}
  \gradx{1,2,3,4,5,6} \grady{-4,-2,2,4,6,8}
  \foreach \p in {(0,8),(1,6),(2,4),(3,2),(4,0)} \path[coulCourbe] \p node[croix]{};
  \foreach \p in {(5,-2),(6,-4)} \path[coulCourbeB] \p node[croixdroite]{};
\end{tikzpicture}

\gris{croix droites : $u_{5}$ et $u_{6}$, dans le prolongement}
\end{minipage}
\end{exo}

\partie{Sens de variation d'une suite}

\begin{exo}[Conjecturer]
\explique{On regarde si chaque terme est plus grand ou plus petit que le précédent, en
calculant les écarts.}{quelques termes donnent une idée, sans prouver}
\cas{a)}
\explique{Écarts $+2$ ; $+3$ ; $+4$ ; $+5$ : les termes augmentent. La suite semble
\res{croissante}.}{chaque terme est plus grand que le précédent}
\cas{b)}
\explique{Écarts $-3$ ; $-4$ ; $-5$ ; $-6$ : les termes diminuent. La suite semble
\res{décroissante}.}{}
\cas{c)}
\explique{Les termes descendent ($2$ puis $-2$), puis remontent ($-2$ puis $2$) : ni
croissante, ni décroissante.}{il faut que tous les écarts aient le même signe}
\explique{La suite ne semble \res{ni croissante, ni décroissante}.}{}
\cas{d)}
\explique{Écarts $+2$ ; $+1$ ; $+2$ ; $+1$ : les termes augmentent toujours. La suite
semble \res{croissante}.}{les écarts changent, mais restent positifs}
\rem{Quelques termes ne prouvent rien : c'est une conjecture. La preuve passe par le signe
de $u_{n+1}-u_{n}$ (exercices suivants).}
\end{exo}

\begin{exo}[Suites définies par récurrence]
\explique{La relation donne $u_{n+1}$ à partir de $u_{n}$ : on enlève $u_{n}$ des deux
côtés, et la différence se lit directement.}{les $u_{n}$ se simplifient}
\cas{a)}
\begin{calculs}
u_{n+1}-u_{n} &= -6 &&\justif{on lit la différence dans la relation}
\end{calculs}
\explique{$u_{n+1}-u_{n}<0$ : la suite est \res{décroissante}.}{différence négative :
chaque terme est plus petit que le précédent}
\cas{b)}
\begin{calculs}
u_{n+1}-u_{n} &= n^{2}+1 &&\justif{on lit la différence dans la relation}
\end{calculs}
\explique{Un carré est positif : $n^{2}\ge 0$, donc $n^{2}+1\ge 1$.}{le signe dépend de
$n$ : on le justifie}
\explique{$u_{n+1}-u_{n}>0$ : la suite est \res{croissante}.}{}
\cas{c)}
\begin{calculs}
u_{n+1}-u_{n} &= -2n^{2} &&\justif{on lit la différence dans la relation}
\end{calculs}
\explique{$n^{2}\ge 0$, donc $-2n^{2}\le 0$.}{multiplier par $-2$ change le sens}
\explique{$u_{n+1}-u_{n}\le 0$ : la suite est \res{décroissante}.}{chaque terme est
inférieur ou égal au précédent ($u_{1}=u_{0}$)}
\cas{d)}
\begin{calculs}
u_{n+1}-u_{n} &= 5n+3 &&\justif{on lit la différence dans la relation}
\end{calculs}
\explique{Comme $n\ge 0$, on a $5n\ge 0$, donc $5n+3\ge 3$.}{$n$ est un rang : il est
positif}
\explique{$u_{n+1}-u_{n}>0$ : la suite est \res{croissante}.}{}
\rem{Le premier terme ne sert pas : le sens de variation ne dépend que du signe de la
différence.}
\end{exo}

\begin{exo}[Suites définies explicitement]
\cas{a)}
\begin{calculs}
u_{n+1} &= 4(n+1)-9 &&\justif{on remplace $n$ par $n+1$}\\
u_{n+1} &= 4n+4-9 &&\justif{on développe}\\
u_{n+1} &= 4n-5 &&
\end{calculs}
\begin{calculs}
u_{n+1}-u_{n} &= 4n-5-(4n-9) &&\justif{parenthèses autour de $u_{n}$}\\
u_{n+1}-u_{n} &= 4n-5-4n+9 &&\justif{le signe $-$ change les signes de la parenthèse}\\
u_{n+1}-u_{n} &= 4 &&
\end{calculs}
\explique{$u_{n+1}-u_{n}>0$ : la suite est \res{croissante}.}{}
\cas{b)}
\begin{calculs}
u_{n+1} &= 6-3(n+1) &&\\
u_{n+1} &= 6-3n-3 &&\\
u_{n+1} &= 3-3n &&
\end{calculs}
\begin{calculs}
u_{n+1}-u_{n} &= 3-3n-(6-3n) &&\\
u_{n+1}-u_{n} &= 3-3n-6+3n &&\\
u_{n+1}-u_{n} &= -3 &&
\end{calculs}
\explique{$u_{n+1}-u_{n}<0$ : la suite est \res{décroissante}.}{}
\cas{c)}
\begin{calculs}
u_{n+1} &= 2(n+1)^{2}-5 &&\\
u_{n+1} &= 2\left(n^{2}+2n+1\right)-5 &&\justif{identité remarquable}\\
u_{n+1} &= 2n^{2}+4n+2-5 &&\\
u_{n+1} &= 2n^{2}+4n-3 &&
\end{calculs}
\begin{calculs}
u_{n+1}-u_{n} &= 2n^{2}+4n-3-\left(2n^{2}-5\right) &&\\
u_{n+1}-u_{n} &= 2n^{2}+4n-3-2n^{2}+5 &&\\
u_{n+1}-u_{n} &= 4n+2 &&
\end{calculs}
\explique{Comme $n\ge 0$, on a $4n+2\ge 2$.}{le signe de la différence dépend de $n$}
\explique{Donc $u_{n+1}-u_{n}>0$ : la suite est \res{croissante}.}{}
\cas{d)}
\begin{calculs}
u_{n+1} &= (n+1)^{2}+6(n+1) &&\\
u_{n+1} &= n^{2}+2n+1+6n+6 &&\justif{identité remarquable}\\
u_{n+1} &= n^{2}+8n+7 &&
\end{calculs}
\begin{calculs}
u_{n+1}-u_{n} &= n^{2}+8n+7-\left(n^{2}+6n\right) &&\\
u_{n+1}-u_{n} &= 2n+7 &&
\end{calculs}
\explique{Comme $n\ge 0$, on a $2n+7\ge 7$.}{}
\explique{Donc $u_{n+1}-u_{n}>0$ : la suite est \res{croissante}.}{}
\rem{Erreur fréquente : oublier les parenthèses autour de $u_{n}$. On soustrait toute
l'expression : $-(6-3n)=-6+3n$.}
\end{exo}

\begin{exo}[Conjecture, puis preuve]
\cas{1)}
\begin{calculs}
u_{0} &= -0^{2}+4\times 0+5 = 5 &&\justif{on remplace $n$ par $0$}\\
u_{1} &= -1^{2}+4\times 1+5 = 8 &&\justif{$-1+4+5$ : le carré, puis le signe $-$}\\
u_{2} &= -2^{2}+4\times 2+5 = 9 &&\justif{$-4+8+5$}\\
u_{3} &= -3^{2}+4\times 3+5 = 8 &&\justif{$-9+12+5$}\\
u_{4} &= -4^{2}+4\times 4+5 = 5 &&\justif{$-16+16+5$}\\
u_{5} &= -5^{2}+4\times 5+5 = 0 &&\justif{$-25+20+5$}
\end{calculs}
\explique{Les termes $5$ ; $8$ ; $9$ ; $8$ ; $5$ ; $0$ montent jusqu'à $u_{2}$, puis
descendent : à partir de $u_{2}$, la suite \res{semble décroissante}.}{$-n^{2}$ : on
calcule $n^{2}$, puis on prend l'opposé}
\cas{2)}
\begin{calculs}
u_{n+1} &= -(n+1)^{2}+4(n+1)+5 &&\justif{on remplace $n$ par $n+1$}\\
u_{n+1} &= -\left(n^{2}+2n+1\right)+4n+4+5 &&\justif{identité remarquable}\\
u_{n+1} &= -n^{2}-2n-1+4n+9 &&\justif{le signe $-$ porte sur les trois termes}\\
u_{n+1} &= -n^{2}+2n+8 &&
\end{calculs}
\begin{calculs}
u_{n+1}-u_{n} &= -n^{2}+2n+8-\left(-n^{2}+4n+5\right) &&\justif{parenthèses autour de $u_{n}$}\\
u_{n+1}-u_{n} &= -n^{2}+2n+8+n^{2}-4n-5 &&\\
u_{n+1}-u_{n} &= \res{3-2n} &&
\end{calculs}
\explique{Contrôle avec $n=0$ : $u_{1}-u_{0}=8-5=3$, et $3-2\times 0=3$.}{}
\cas{3)}
\explique{Pour $n\ge 2$, on a $2n\ge 4$, donc $3-2n\le 3-4$, soit $3-2n\le -1$.}{on part
de $n\ge 2$}
\explique{$u_{n+1}-u_{n}<0$ dès que $n\ge 2$ : la suite \res{est décroissante à partir du
rang $2$}.}{la conjecture est prouvée}
\rem{Pour $n=0$, $3-2n$ vaut $3$ ; pour $n=1$, il vaut $1$ : la suite monte d'abord, comme on l'a
vu à la question 1).}
\end{exo}

\partie{Suites arithmétiques}

\begin{exo}[Premiers termes]
\explique{Une suite arithmétique ajoute toujours la raison : $u_{n+1}=u_{n}+r$.}{on
calcule chaque terme à partir du précédent}
\cas{a)}
\explique{$u_{n+1}=\res{u_{n}+6}$}{on ajoute la raison $6$}
\begin{calculs}
u_{1} &= 4+6 = \res{10} &&\\
u_{2} &= 10+6 = \res{16} &&\\
u_{3} &= 16+6 = \res{22} &&
\end{calculs}
\cas{b)}
\explique{$u_{n+1}=\res{u_{n}-5}$}{ajouter $-5$, c'est enlever $5$}
\begin{calculs}
u_{1} &= 12-5 = \res{7} &&\\
u_{2} &= 7-5 = \res{2} &&\\
u_{3} &= 2-5 = \res{-3} &&
\end{calculs}
\cas{c)}
\explique{$u_{n+1}=\res{u_{n}+0{,}5}$}{}
\begin{calculs}
u_{1} &= -3+0{,}5 = \res{-2{,}5} &&\\
u_{2} &= -2{,}5+0{,}5 = \res{-2} &&\\
u_{3} &= -2+0{,}5 = \res{-1{,}5} &&
\end{calculs}
\cas{d)}
\explique{$u_{n+1}=\res{u_{n}+\dfrac{1}{2}}$, et
$\dfrac{1}{2}=\dfrac{1\fois{2}}{2\fois{2}}=\dfrac{2}{4}$.}{on met la raison au même
dénominateur que $u_{0}$}
\begin{calculs}
u_{1} &= \dfrac{1}{4}+\dfrac{2}{4} = \res{\dfrac{3}{4}} &&\justif{on ajoute les numérateurs}\\
u_{2} &= \dfrac{3}{4}+\dfrac{2}{4} = \res{\dfrac{5}{4}} &&\\
u_{3} &= \dfrac{5}{4}+\dfrac{2}{4} = \res{\dfrac{7}{4}} &&
\end{calculs}
\end{exo}

\begin{exo}[Arithmétique ou pas ?]
\explique{Une situation se traduit par une suite arithmétique quand on \textbf{ajoute} (ou
enlève) toujours le \textbf{même nombre}.}{un pourcentage ou une fraction ne donne pas
une suite arithmétique}
\cas{a)}
\explique{Chaque année, on ajoute $15$~€ : \res{oui}, de premier terme \res{$u_{0}=650$}
et de raison \res{$r=15$}.}{$u_{n+1}=u_{n}+15$}
\cas{b)}
\explique{$+10\,\%$ : $5\,000$ ; $5\,500$ ; $6\,050$. Les écarts $500$ et $550$ ne sont
pas égaux : \res{non}.}{$10\,\%$ de $5\,500$ vaut $550$ : l'augmentation grandit}
\cas{c)}
\explique{Chaque heure, on enlève $4~\text{m}^{3}$ : \res{oui}, de premier terme
\res{$u_{0}=60$} et de raison \res{$r=-4$}.}{la raison est négative quand on enlève}
\cas{d)}
\explique{$3$ ; $2{,}25$ ; $1{,}687\,5$ : les écarts $-0{,}75$ et $-0{,}562\,5$ ne sont
pas égaux : \res{non}.}{les trois quarts : on multiplie par $0{,}75$}
\rem{« On ajoute le même nombre » : arithmétique. « On augmente du même pourcentage » ou
« on prend la même fraction » : ce n'est pas arithmétique.}
\end{exo}

\begin{exo}[Le food truck]
\cas{1)}
\explique{Chaque semaine, le food truck sert $5$ repas de plus que la semaine
précédente : $u_{n+1}=u_{n}+5$.}{on ajoute toujours le même nombre}
\explique{La suite $(u_{n})$ est \res{arithmétique, de raison $5$}.}{premier terme :
$u_{1}=40$}
\cas{2)}
\begin{calculs}
u_{2} &= 40+5 = \res{45} &&\justif{la $2^{\text{e}}$ semaine : $u_{2}$}\\
u_{3} &= 45+5 = 50 &&\\
u_{4} &= 50+5 = 55 &&\\
u_{5} &= 55+5 = 60 &&\\
u_{6} &= 60+5 = \res{65} &&\justif{la $6^{\text{e}}$ semaine : $u_{6}$}
\end{calculs}
\explique{Le food truck sert $45$ repas la $2^{\text{e}}$ semaine et $65$ la
$6^{\text{e}}$.}{ici, la suite commence à $u_{1}$ : $u_{6}$ est bien la $6^{\text{e}}$
semaine}
\end{exo}

\begin{exo}[Une application]
\cas{1)}
\begin{calculs}
v_{1} &= \res{2\,500} &&\justif{le premier mois}\\
v_{2} &= 2\,500+400 = \res{2\,900} &&\justif{$400$ utilisateurs de plus}\\
v_{3} &= 2\,900+400 = \res{3\,300} &&\\
v_{4} &= 3\,300+400 = \res{3\,700} &&
\end{calculs}
\cas{2)}
\explique{Chaque mois, on ajoute $400$ : $v_{n+1}=v_{n}+400$.}{toujours le même nombre}
\explique{La suite $(v_{n})$ est \res{arithmétique, de raison $400$}.}{premier terme :
$v_{1}=2\,500$}
\end{exo}

\partie{Reconnaître une suite arithmétique}

\begin{exo}[À partir des termes]
\explique{On calcule \textbf{toutes} les différences entre deux termes qui se
suivent.}{deux différences distinctes suffisent pour dire non}
\cas{a)}
\explique{$8-3=5$ ; \ $13-8=5$ ; \ $18-13=5$ : \res{oui, de raison $5$}.}{}
\cas{b)}
\explique{$14-20=-6$ ; \ $8-14=-6$ ; \ $2-8=-6$ : \res{oui, de raison $-6$}.}{la liste
descend : raison négative}
\cas{c)}
\explique{$3-2=1$ mais $5-3=2$ : \res{non}.}{}
\cas{d)}
\explique{$2{,}25-1{,}5=0{,}75$ ; \ $3-2{,}25=0{,}75$ ; \ $3{,}75-3=0{,}75$ : \res{oui, de
raison $0{,}75$}.}{}
\cas{e)}
\explique{$377-401=-24$ ; \ $353-377=-24$ ; \ $329-353=-24$ : \res{oui, de raison
$-24$}.}{}
\cas{f)}
\explique{$4-2=2$ mais $8-4=4$ : \res{non}.}{on multiplie par $2$ : ce n'est pas une
addition}
\end{exo}

\begin{exo}[Démontrer]
\cas{1)}
\begin{calculs}
u_{0} &= 5\times 0-2 = \res{-2} &&\\
u_{1} &= 5\times 1-2 = \res{3} &&\\
u_{2} &= 5\times 2-2 = \res{8} &&\\
u_{3} &= 5\times 3-2 = \res{13} &&\justif{on ajoute $5$ à chaque fois}
\end{calculs}
\cas{2)}
\begin{calculs}
u_{n+1} &= 5(n+1)-2 &&\justif{on remplace $n$ par $n+1$}\\
u_{n+1} &= 5n+5-2 &&\justif{on développe}\\
u_{n+1} &= \res{5n+3} &&
\end{calculs}
\cas{3)}
\begin{calculs}
u_{n+1}-u_{n} &= 5n+3-(5n-2) &&\justif{parenthèses autour de $u_{n}$}\\
u_{n+1}-u_{n} &= 5n+3-5n+2 &&\\
u_{n+1}-u_{n} &= 5 &&\justif{un nombre, sans $n$}
\end{calculs}
\explique{\res{$(u_{n})$ est arithmétique, de premier terme $u_{0}=-2$ et de raison
$r=5$}.}{la différence est la même pour tout $n$}
\rem{Les termes de la question 1) donnent une idée, mais ne démontrent rien : il faut le
calcul de $u_{n+1}-u_{n}$ pour tout $n$.}
\end{exo}

\begin{exo}[Pour chaque suite]
\explique{Pour montrer qu'une suite est arithmétique, on calcule $u_{n+1}-u_{n}$ ; pour
montrer qu'elle ne l'est pas, deux différences distinctes entre les premiers termes
suffisent.}{méthode 8 du cours}
\cas{a)}
\begin{calculs}
a_{n+1} &= 4(n+1) &&\justif{on remplace $n$ par $n+1$}\\
a_{n+1} &= 4n+4 &&\\
a_{n+1}-a_{n} &= 4n+4-4n &&\\
a_{n+1}-a_{n} &= 4 &&\justif{un nombre, sans $n$}
\end{calculs}
\explique{\res{Arithmétique, de raison $r=4$ et de premier terme $a_{0}=0$}.}{$a_{0}=4\times
0$}
\cas{b)}
\begin{calculs}
b_{n+1} &= 2-3(n+1) &&\\
b_{n+1} &= -1-3n &&\justif{$2-3n-3$}\\
b_{n+1}-b_{n} &= -1-3n-(2-3n) &&\\
b_{n+1}-b_{n} &= -3 &&\justif{$-1-3n-2+3n$}
\end{calculs}
\explique{\res{Arithmétique, de raison $r=-3$ et de premier terme $b_{0}=2$}.}{}
\cas{c)}
\explique{$c_{0}=0$ ; $c_{1}=2$ ; $c_{2}=6$ ; or $c_{1}-c_{0}=2$ et $c_{2}-c_{1}=4$.}{les
premiers termes suffisent pour dire non}
\explique{Les différences ne sont pas égales : \res{$(c_{n})$ n'est pas
arithmétique}.}{}
\cas{d)}
\begin{calculs}
d_{n+1}-d_{n} &= \dfrac{n+1}{3}-\dfrac{n}{3} &&\justif{$d_{n+1}=\dfrac{n+1}{3}$}\\
d_{n+1}-d_{n} &= \dfrac{1}{3} &&\justif{même dénominateur : $n+1-n=1$}
\end{calculs}
\explique{\res{Arithmétique, de raison $r=\dfrac{1}{3}$ et de premier terme
$d_{0}=0$}.}{}
\cas{e)}
\explique{$e_{0}=2$ ; $e_{1}=1$ ; $e_{2}=-2$ ; or $e_{1}-e_{0}=-1$ et
$e_{2}-e_{1}=-3$.}{$e_{2}=2-4$}
\explique{\res{$(e_{n})$ n'est pas arithmétique}.}{}
\cas{f)}
\explique{On développe d'abord $f_{n}$ :}{l'expression se simplifie}
\begin{calculs}
f_{n} &= (n+2)^{2}-n^{2} &&\\
f_{n} &= n^{2}+4n+4-n^{2} &&\justif{identité remarquable}\\
f_{n} &= 4n+4 &&\\
f_{n+1}-f_{n} &= 4(n+1)+4-(4n+4) &&\\
f_{n+1}-f_{n} &= 4 &&\justif{$4n+8-4n-4$}
\end{calculs}
\explique{\res{Arithmétique, de raison $r=4$ et de premier terme $f_{0}=4$}.}{$f_{0}=2^{2}-0^{2}$}
\cas{g)}
\begin{calculs}
g_{1} &= 1^{2}+1 = 2 &&\justif{$g_{1}=g_{0}^{2}+1$}\\
g_{2} &= 2^{2}+1 = 5 &&
\end{calculs}
\explique{$g_{1}-g_{0}=1$ et $g_{2}-g_{1}=3$ : \res{$(g_{n})$ n'est pas
arithmétique}.}{deux différences distinctes}
\end{exo}

\partie{Terme général et variations d'une suite arithmétique}

\begin{exo}[Calculer un terme éloigné]
\explique{On applique $u_{n}=u_{0}+n\times r$, puis le signe de $r$ donne le sens de
variation.}{propriétés 2 et 3}
\cas{a)}
\begin{calculs}
u_{n} &= 8+n\times(-3) &&\justif{$u_{n}=u_{0}+n\times r$}\\
u_{n} &= \res{8-3n} &&\\
u_{9} &= 8-3\times 9 &&\justif{on remplace $n$ par $9$}\\
u_{9} &= 8-27 = \res{-19} &&
\end{calculs}
\explique{$r=-3<0$ : la suite est \res{décroissante}.}{signe de la raison}
\cas{b)}
\begin{calculs}
u_{n} &= \res{2+6n} &&\justif{$u_{n}=u_{0}+n\times r$}\\
u_{100} &= 2+6\times 100 &&\\
u_{100} &= 2+600 = \res{602} &&
\end{calculs}
\explique{$r=6>0$ : la suite est \res{croissante}.}{}
\cas{c)}
\begin{calculs}
u_{n} &= \res{-5+\dfrac{5}{2}n} &&\\
u_{6} &= -5+\dfrac{5}{2}\times 6 &&\\
u_{6} &= -5+15 = \res{10} &&\justif{$\dfrac{5\times 6}{2}=\dfrac{30}{2}$}
\end{calculs}
\explique{$r=\dfrac{5}{2}>0$ : la suite est \res{croissante}.}{}
\cas{d)}
\begin{calculs}
u_{n} &= 4+n\times\left(-\dfrac{2}{3}\right) &&\\
u_{n} &= \res{4-\dfrac{2}{3}n} &&\\
u_{9} &= 4-\dfrac{2}{3}\times 9 &&\\
u_{9} &= 4-6 = \res{-2} &&\justif{$\dfrac{2\times 9}{3}=\dfrac{18}{3}$}
\end{calculs}
\explique{$r=-\dfrac{2}{3}<0$ : la suite est \res{décroissante}.}{}
\end{exo}

\begin{exo}[Quand on ne part pas de $u_{0}$]
\cas{a)}
\begin{calculs}
u_{n} &= u_{1}+(n-1)\times r &&\justif{la suite commence à $u_{1}$}\\
u_{n} &= 3+(n-1)\times 4 &&\\
u_{n} &= 3+4n-4 &&\justif{on développe}\\
u_{n} &= \res{4n-1} &&\\
u_{10} &= 4\times 10-1 = \res{39} &&
\end{calculs}
\explique{Contrôle : $u_{1}=4\times 1-1=3$.}{on retrouve le premier terme}
\cas{b)}
\begin{calculs}
u_{n} &= -6+(n-1)\times 0{,}5 &&\justif{$u_{n}=u_{1}+(n-1)\times r$}\\
u_{21} &= -6+20\times 0{,}5 &&\justif{$n-1=20$}\\
u_{21} &= -6+10 = \res{4} &&
\end{calculs}
\rem{De $u_{1}$ à $u_{21}$, on ajoute $20$ fois la raison, pas $21$ fois.}
\cas{c)}
\explique{De $u_{0}$ à $u_{4}$, on ajoute $4$ fois la raison : $u_{4}=u_{0}+4\times
3$.}{$u_{n}=u_{0}+n\times r$ avec $n=4$}
\begin{calculs}
u_{0}+4\times 3 &= 10 &&\justif{on traduit $u_{4}=10$}\\
u_{0}+12 &= 10 &&\\
u_{0} &= \res{-2} &&\justif{on enlève $12$ des deux côtés}
\end{calculs}
\begin{calculs}
u_{12} &= u_{0}+12\times r &&\\
u_{12} &= -2+12\times 3 &&\\
u_{12} &= -2+36 = \res{34} &&
\end{calculs}
\explique{Contrôle : de $u_{4}$ à $u_{12}$, on ajoute $8$ fois $3$ ; $10+24=34$.}{}
\end{exo}

\begin{exo}[Une application (suite)]
\cas{1)}
\begin{calculs}
v_{n} &= v_{1}+(n-1)\times r &&\justif{la suite commence à $v_{1}$}\\
v_{n} &= \res{2\,500+(n-1)\times 400} &&
\end{calculs}
\cas{2)}
\begin{calculs}
v_{12} &= 2\,500+11\times 400 &&\justif{le $12^{\text{e}}$ mois : $n=12$, donc $n-1=11$}\\
v_{12} &= 2\,500+4\,400 = \res{6\,900} &&
\end{calculs}
\explique{Le $12^{\text{e}}$ mois, l'application comptera $6\,900$ utilisateurs.}{une
phrase pour conclure}
\cas{3)}
\explique{$r=400>0$ : \res{$(v_{n})$ est croissante}.}{signe de la raison}
\end{exo}

\begin{exo}[Économiser pour un scooter]
\cas{1)}
\begin{calculs}
u_{2} &= 1\,100+75 = \res{1\,175} &&\justif{au $1^{\text{er}}$ avril}\\
u_{3} &= 1\,175+75 = \res{1\,250} &&\justif{au $1^{\text{er}}$ mai}\\
u_{4} &= 1\,250+75 = \res{1\,325} &&\justif{au $1^{\text{er}}$ juin}
\end{calculs}
\cas{2)}
\explique{On ajoute $75$ chaque mois : $(u_{n})$ est arithmétique, de raison $75$ et de
premier terme $u_{1}=1\,100$.}{}
\begin{calculs}
u_{n} &= u_{1}+(n-1)\times r &&\justif{la suite commence à $u_{1}$}\\
u_{n} &= \res{1\,100+(n-1)\times 75} &&\\
u_{7} &= 1\,100+6\times 75 &&\\
u_{7} &= 1\,100+450 = \res{1\,550} &&
\end{calculs}
\cas{3)}
\begin{calculs}
1\,100+(n-1)\times 75 &= 2\,000 &&\justif{on traduit $u_{n}=2\,000$}\\
(n-1)\times 75 &= 900 &&\justif{on enlève $1\,100$ des deux côtés}\\
n-1 &= 12 &&\justif{on divise par $75$}\\
n &= 13 &&
\end{calculs}
\explique{$u_{13}$ est l'épargne $12$ mois après le $1^{\text{er}}$ mars : elle pourra
acheter le scooter \res{le $1^{\text{er}}$ mars de l'année suivante}.}{contrôle :
$1\,100+12\times 75=1\,100+900=2\,000$}
\end{exo}

\begin{exo}[La production d'une usine]
\cas{1)}
\begin{calculs}
p_{n} &= p_{1}+(n-1)\times r &&\justif{la suite commence à $p_{1}$}\\
p_{n} &= \res{4\,500+(n-1)\times 500} &&
\end{calculs}
\cas{2)}
\explique{$p_{1}$ est 2024 : l'année $2023+n$ correspond au rang $n$. 2027 donne $n=4$,
2040 donne $n=17$.}{on enlève 2023 à l'année}
\begin{calculs}
p_{4} &= 4\,500+3\times 500 &&\\
p_{4} &= 4\,500+1\,500 = \res{6\,000} &&\\
p_{17} &= 4\,500+16\times 500 &&\\
p_{17} &= 4\,500+8\,000 = \res{12\,500} &&
\end{calculs}
\explique{L'usine fabriquera $6\,000$ pièces en 2027 et $12\,500$ pièces en 2040.}{}
\cas{3)}
\explique{Tripler, c'est atteindre $3\times 4\,500=13\,500$ pièces.}{on cherche $n$ tel que
$p_{n}=13\,500$}
\begin{calculs}
4\,500+(n-1)\times 500 &= 13\,500 &&\\
(n-1)\times 500 &= 9\,000 &&\justif{on enlève $4\,500$}\\
n-1 &= 18 &&\justif{on divise par $500$}\\
n &= 19 &&
\end{calculs}
\explique{$p_{19}$ correspond à l'année $2023+19$ : la production aura triplé \res{en
2042}.}{}
\rem{Erreur fréquente : $2024+19=2043$. Comme $p_{1}$ est 2024, c'est $2023+n$.}
\end{exo}

\partie{Suites géométriques}

\begin{exo}[Coefficient multiplicateur]
\explique{Augmenter de $t\,\%$, c'est multiplier par $1+\dfrac{t}{100}$ ; diminuer de
$t\,\%$, par $1-\dfrac{t}{100}$.}{propriété 4}
\cas{a)}
\explique{$q=1+0{,}04=\res{1{,}04}$}{$4\,\%=\dfrac{4}{100}=0{,}04$}
\cas{b)}
\explique{$q=1-0{,}04=\res{0{,}96}$}{}
\cas{c)}
\explique{$q=1+0{,}3=\res{1{,}3}$}{$30\,\%=0{,}30$}
\cas{d)}
\explique{$q=1-0{,}3=\res{0{,}7}$}{}
\cas{e)}
\explique{$q=1+0{,}65=\res{1{,}65}$}{}
\cas{f)}
\explique{$q=1-0{,}12=\res{0{,}88}$}{}
\cas{g)}
\explique{$q=1+2=\res{3}$}{$200\,\%=\dfrac{200}{100}=2$}
\cas{h)}
\explique{$q=1-0{,}45=\res{0{,}55}$}{}
\rem{Augmenter de $200\,\%$, c'est tripler (multiplier par $3$), et non doubler.}
\end{exo}

\begin{exo}[Dans l'autre sens]
\explique{Si $q>1$, c'est une augmentation de $q-1$ ; si $q<1$, une diminution de $1-q$.
On écrit ensuite ce nombre en pourcentage.}{$0{,}04=4\,\%$}
\cas{a)}
\explique{$0{,}96<1$ et $1-0{,}96=0{,}04$ : \res{diminution de $4\,\%$}.}{}
\cas{b)}
\explique{$1{,}07>1$ et $1{,}07-1=0{,}07$ : \res{augmentation de $7\,\%$}.}{}
\cas{c)}
\explique{$0{,}25<1$ et $1-0{,}25=0{,}75$ : \res{diminution de $75\,\%$}.}{}
\cas{d)}
\explique{$1{,}6>1$ et $1{,}6-1=0{,}6$ : \res{augmentation de $60\,\%$}.}{$0{,}6=0{,}60$}
\cas{e)}
\explique{$3>1$ et $3-1=2$ : \res{augmentation de $200\,\%$}.}{$2=\dfrac{200}{100}$}
\cas{f)}
\explique{$0{,}98<1$ et $1-0{,}98=0{,}02$ : \res{diminution de $2\,\%$}.}{}
\cas{g)}
\explique{$1{,}035>1$ et $1{,}035-1=0{,}035$ : \res{augmentation de
$3{,}5\,\%$}.}{$0{,}035=\dfrac{3{,}5}{100}$}
\cas{h)}
\explique{$0{,}82<1$ et $1-0{,}82=0{,}18$ : \res{diminution de $18\,\%$}.}{}
\end{exo}

\begin{exo}[Premiers termes]
\explique{Une suite géométrique multiplie toujours par la raison : $u_{n+1}=q\times
u_{n}$.}{chaque terme à partir du précédent}
\cas{a)}
\begin{calculs}
u_{1} &= 3\times 5 = \res{15} &&\justif{$u_{1}=q\times u_{0}$}\\
u_{2} &= 3\times 15 = \res{45} &&\\
u_{3} &= 3\times 45 = \res{135} &&
\end{calculs}
\cas{b)}
\begin{calculs}
u_{1} &= 0{,}5\times 96 = \res{48} &&\justif{multiplier par $0{,}5$ : prendre la moitié}\\
u_{2} &= 0{,}5\times 48 = \res{24} &&\\
u_{3} &= 0{,}5\times 24 = \res{12} &&
\end{calculs}
\cas{c)}
\begin{calculs}
u_{1} &= 1{,}05\times 2\,000 = \res{2\,100} &&\justif{$+5\,\%$}\\
u_{2} &= 1{,}05\times 2\,100 = \res{2\,205} &&\\
u_{3} &= 1{,}05\times 2\,205 = \res{2\,315{,}25} &&
\end{calculs}
\cas{d)}
\begin{calculs}
u_{1} &= \dfrac{1}{3}\times 81 = \res{27} &&\justif{multiplier par $\dfrac{1}{3}$ : diviser
par $3$}\\
u_{2} &= \dfrac{1}{3}\times 27 = \res{9} &&\\
u_{3} &= \dfrac{1}{3}\times 9 = \res{3} &&
\end{calculs}
\end{exo}

\begin{exo}[Deux situations]
\cas{a)}
\explique{$u_{0}=\res{3\,000}$ et $u_{n+1}=\res{1{,}02\times u_{n}}$.}{$+2\,\%$ : on
multiplie par $1+0{,}02$}
\begin{calculs}
u_{1} &= 1{,}02\times 3\,000 = \res{3\,060} &&\\
u_{2} &= 1{,}02\times 3\,060 = \res{3\,121{,}20} &&\justif{les intérêts portent aussi sur les
intérêts}
\end{calculs}
\explique{Le capital vaut $3\,060$~€ après un an, $3\,121{,}20$~€ après deux ans.}{}
\cas{b)}
\explique{$u_{0}=\res{40\,000}$ et $u_{n+1}=\res{0{,}97\times u_{n}}$.}{$-3\,\%$ : on
multiplie par $1-0{,}03$}
\begin{calculs}
u_{1} &= 0{,}97\times 40\,000 = \res{38\,800} &&\\
u_{2} &= 0{,}97\times 38\,800 = \res{37\,636} &&
\end{calculs}
\explique{La ville compte $38\,800$ habitants après un an, $37\,636$ après deux
ans.}{}
\rem{Erreur fréquente : enlever $3\,\%$ de $40\,000$ chaque année ($1\,200$ habitants).
Les $3\,\%$ portent sur la population de l'année précédente.}
\end{exo}

\partie{Reconnaître une suite géométrique}

\begin{exo}[À partir des termes]
\explique{On calcule \textbf{tous} les quotients d'un terme par le précédent.}{deux
quotients distincts suffisent pour dire non}
\cas{a)}
\explique{$\dfrac{10}{5}=2$ ; \ $\dfrac{20}{10}=2$ ; \ $\dfrac{40}{20}=2$ : \res{oui, de
raison $2$}.}{}
\cas{b)}
\explique{$\dfrac{12}{3}=4$ ; \ $\dfrac{48}{12}=4$ ; \ mais $\dfrac{180}{48}=3{,}75$ :
\res{non}.}{$48\times 4=192$, pas $180$}
\cas{c)}
\explique{$\dfrac{32}{128}=\dfrac{1}{4}$ ; \ $\dfrac{8}{32}=\dfrac{1}{4}$ ; \
$\dfrac{2}{8}=\dfrac{1}{4}$ : \res{oui, de raison $\dfrac{1}{4}$}.}{on divise par $4$ à
chaque fois}
\cas{d)}
\explique{$\dfrac{8}{4}=2$ mais $\dfrac{12}{8}=1{,}5$ : \res{non}.}{c'est une suite
arithmétique, de raison $4$}
\cas{e)}
\explique{$\dfrac{220}{200}=1{,}1$ ; \ $\dfrac{242}{220}=1{,}1$ ; \
$\dfrac{266{,}2}{242}=1{,}1$ : \res{oui, de raison $1{,}1$}.}{une hausse de $10\,\%$ à
chaque fois}
\cas{f)}
\explique{$\dfrac{10}{50}=0{,}2$ ; \ $\dfrac{2}{10}=0{,}2$ ; \ $\dfrac{0{,}4}{2}=0{,}2$ :
\res{oui, de raison $0{,}2$}.}{}
\rem{En b), les deux premiers quotients sont égaux : il faut aller jusqu'au bout de la
liste.}
\end{exo}

\begin{exo}[Démontrer]
\cas{1)}
\begin{calculs}
u_{0} &= 2\times 3^{0} = \res{2} &&\justif{$3^{0}=1$}\\
u_{1} &= 2\times 3^{1} = \res{6} &&\\
u_{2} &= 2\times 3^{2} = \res{18} &&\justif{la puissance d'abord : $2\times 9$}\\
u_{3} &= 2\times 3^{3} = \res{54} &&\justif{$2\times 27$}
\end{calculs}
\cas{2)}
\begin{calculs}
u_{n+1} &= \res{2\times 3^{n+1}} &&\justif{on remplace $n$ par $n+1$}
\end{calculs}
\cas{3)}
\begin{calculs}
u_{n+1} &= 2\times 3^{n+1} &&\\
u_{n+1} &= 2\times 3^{n}\times 3 &&\justif{$3^{n+1}=3^{n}\times 3$}\\
u_{n+1} &= 3\times u_{n} &&\justif{on reconnaît $u_{n}=2\times 3^{n}$}
\end{calculs}
\explique{\res{$(u_{n})$ est géométrique, de premier terme $u_{0}=2$ et de raison
$q=3$}.}{on multiplie toujours par $3$}
\end{exo}

\begin{exo}[Pour chaque suite]
\explique{Pour montrer qu'une suite est géométrique, on écrit $u_{n+1}=q\times u_{n}$ ;
pour montrer qu'elle ne l'est pas, deux quotients distincts entre les premiers termes
suffisent.}{méthode 11 du cours}
\cas{a)}
\begin{calculs}
a_{n+1} &= 5^{n+1} &&\justif{on remplace $n$ par $n+1$}\\
a_{n+1} &= 5^{n}\times 5 &&\\
a_{n+1} &= 5\times a_{n} &&
\end{calculs}
\explique{\res{Géométrique, de raison $q=5$ et de premier terme $a_{0}=1$}.}{$5^{0}=1$}
\cas{b)}
\explique{$b_{0}=3$ ; $b_{1}=4$ ; $b_{2}=7$ ; or $\dfrac{b_{1}}{b_{0}}=\dfrac{4}{3}$ et
$\dfrac{b_{2}}{b_{1}}=\dfrac{7}{4}$.}{les premiers termes suffisent pour dire non}
\explique{Les quotients ne sont pas égaux : \res{$(b_{n})$ n'est pas géométrique}.}{}
\cas{c)}
\begin{calculs}
c_{n+1} &= 8\times 0{,}5^{n+1} &&\\
c_{n+1} &= 8\times 0{,}5^{n}\times 0{,}5 &&\\
c_{n+1} &= 0{,}5\times c_{n} &&
\end{calculs}
\explique{\res{Géométrique, de raison $q=0{,}5$ et de premier terme $c_{0}=8$}.}{$0{,}5^{0}=1$}
\cas{d)}
\explique{$d_{0}=1$ ; $d_{1}=4$ ; $d_{2}=7$ ; or $\dfrac{d_{1}}{d_{0}}=4$ et
$\dfrac{d_{2}}{d_{1}}=\dfrac{7}{4}$.}{}
\explique{\res{$(d_{n})$ n'est pas géométrique}.}{elle est arithmétique, de raison $3$}
\cas{e)}
\begin{calculs}
e_{1} &= 1^{2}+2 = 3 &&\justif{$e_{1}=e_{0}^{2}+2$}\\
e_{2} &= 3^{2}+2 = 11 &&
\end{calculs}
\explique{$\dfrac{e_{1}}{e_{0}}=3$ et $\dfrac{e_{2}}{e_{1}}=\dfrac{11}{3}$ : \res{$(e_{n})$
n'est pas géométrique}.}{deux quotients distincts}
\end{exo}

\begin{exo}[La valeur d'une machine]
\cas{1)}
\explique{$\dfrac{10\,800}{12\,000}=0{,}9$ ; \ $\dfrac{9\,720}{10\,800}=0{,}9$ ; \
$\dfrac{8\,748}{9\,720}=0{,}9$.}{tous les quotients}
\explique{Les quotients sont égaux : ce sont les termes d'une \res{suite géométrique de
raison $0{,}9$}.}{on multiplie par $0{,}9$ chaque année}
\cas{2)}
\explique{$0{,}9<1$ et $1-0{,}9=0{,}1$ : la machine perd \res{$10\,\%$} de sa valeur chaque
année.}{$0{,}1=10\,\%$}
\cas{3)}
\begin{calculs}
v_{4} &= 0{,}9\times v_{3} &&\justif{$v_{n+1}=q\times v_{n}$}\\
v_{4} &= 0{,}9\times 8\,748 = \res{7\,873{,}20} &&
\end{calculs}
\explique{Après $4$ ans, la machine vaut $7\,873{,}20$~€.}{}
\end{exo}

\partie{Terme général et variations d'une suite géométrique}

\begin{exo}[Calculer un terme éloigné]
\explique{On applique $u_{n}=u_{0}\times q^{n}$, puis on compare $q$ à $1$ pour le sens de
variation.}{propriétés 6 et 7 ; ici $u_{0}>0$}
\cas{a)}
\begin{calculs}
u_{n} &= \res{3\times 2^{n}} &&\justif{$u_{n}=u_{0}\times q^{n}$}\\
u_{10} &= 3\times 2^{10} &&\\
u_{10} &= 3\times 1\,024 = \res{3\,072} &&\justif{$2^{10}=1\,024$}
\end{calculs}
\explique{$q=2>1$ : la suite est \res{croissante}.}{}
\cas{b)}
\begin{calculs}
u_{n} &= \res{6\times\left(\dfrac{1}{2}\right)^{n}} &&\\
u_{5} &= 6\times\left(\dfrac{1}{2}\right)^{5} &&\\
u_{5} &= 6\times\dfrac{1}{32} = \dfrac{6}{32} &&\justif{$2^{5}=32$}\\
u_{5} &= \dfrac{3\times\barre{2}}{16\times\barre{2}} = \res{\dfrac{3}{16}} &&\justif{on
simplifie par $2$}
\end{calculs}
\explique{$0<q<1$ : la suite est \res{décroissante}.}{}
\cas{c)}
\begin{calculs}
u_{n} &= \res{2\,000\times 0{,}8^{n}} &&\\
u_{5} &= 2\,000\times 0{,}8^{5} &&\justif{à la calculatrice}\\
u_{5} &= \res{655{,}36} &&
\end{calculs}
\explique{$0<q<1$ : la suite est \res{décroissante}.}{}
\cas{d)}
\begin{calculs}
u_{n} &= 8\times 1^{n} &&\\
u_{n} &= \res{8} &&\justif{$1^{n}=1$}\\
u_{30} &= \res{8} &&
\end{calculs}
\explique{$q=1$ : la suite est \res{constante}.}{multiplier par $1$ ne change rien}
\end{exo}

\begin{exo}[Quand on ne part pas de $u_{0}$]
\cas{a)}
\begin{calculs}
u_{n} &= u_{1}\times q^{n-1} &&\justif{la suite commence à $u_{1}$}\\
u_{n} &= \res{4\times 3^{n-1}} &&\\
u_{6} &= 4\times 3^{5} &&\justif{$n-1=5$}\\
u_{6} &= 4\times 243 = \res{972} &&\justif{$3^{5}=243$}
\end{calculs}
\rem{De $u_{1}$ à $u_{6}$, on multiplie $5$ fois par $q$, pas $6$ fois : l'exposant est
$n-1$.}
\cas{b)}
\begin{calculs}
u_{3} &= 2\times u_{2} &&\justif{$u_{n+1}=q\times u_{n}$}\\
u_{3} &= 2\times 5 = \res{10} &&\\
u_{4} &= 2\times 10 = \res{20} &&\\
u_{5} &= 2\times 20 = \res{40} &&
\end{calculs}
\end{exo}

\begin{exo}[Un placement]
\cas{1)}
\explique{Chaque année, le capital augmente de $5\,\%$ : on le multiplie par
$1{,}05$.}{$1+0{,}05$}
\begin{calculs}
C_{1} &= 1{,}05\times 5\,000 = \res{5\,250} &&\\
C_{2} &= 1{,}05\times 5\,250 = \res{5\,512{,}50} &&
\end{calculs}
\cas{2)}
\explique{$C_{n+1}=1{,}05\times C_{n}$ : $(C_{n})$ est \res{géométrique, de raison
$1{,}05$} et de premier terme $C_{0}=5\,000$.}{on multiplie toujours par $1{,}05$}
\begin{calculs}
C_{n} &= C_{0}\times q^{n} &&\\
C_{n} &= \res{5\,000\times 1{,}05^{n}} &&
\end{calculs}
\cas{3)}
\begin{calculs}
C_{10} &= 5\,000\times 1{,}05^{10} &&\justif{à la calculatrice}\\
C_{10} &\approx \res{8\,144{,}47} &&\justif{arrondi au centime}
\end{calculs}
\explique{Après $10$ ans, le capital vaut environ $8\,144{,}47$~€.}{}
\cas{4)}
\explique{On cherche le premier $n$ tel que $C_{n}\ge 10\,000$, dans le tableau de valeurs
de $5\,000\times 1{,}05^{x}$.}{doubler : atteindre $2\times 5\,000$}
\explique{$C_{14}\approx 9\,899{,}66<10\,000$ \ et \ $C_{15}\approx 10\,394{,}64\ge
10\,000$.}{on lit les deux lignes autour du seuil}
\explique{Le capital aura doublé \res{au bout de $15$ ans}.}{}
\end{exo}

\begin{exo}[La ville (suite)]
\cas{1)}
\begin{calculs}
u_{n} &= u_{0}\times q^{n} &&\\
u_{n} &= \res{40\,000\times 0{,}97^{n}} &&
\end{calculs}
\explique{$0<q<1$ : la suite est \res{décroissante}.}{la population diminue}
\cas{2)}
\begin{calculs}
u_{10} &= 40\,000\times 0{,}97^{10} &&\justif{à la calculatrice}\\
u_{10} &\approx 29\,496{,}97 &&
\end{calculs}
\explique{Après $10$ ans, la ville comptera environ \res{$29\,497$ habitants}.}{une
population s'arrondit à l'unité}
\cas{3)}
\explique{On cherche le premier $n$ tel que $u_{n}<35\,000$, dans le tableau de valeurs de
$40\,000\times 0{,}97^{x}$.}{}
\explique{$u_{4}\approx 35\,412$ \ et \ $u_{5}\approx 34\,349<35\,000$.}{on lit les deux
lignes autour du seuil}
\explique{La ville comptera moins de $35\,000$ habitants \res{au bout de $5$ ans}.}{}
\end{exo}

\begin{exo}[La batterie d'un drone]
\cas{1)}
\explique{À chaque vol, la durée est multipliée par $0{,}98$ : $d_{n+1}=0{,}98\times
d_{n}$. La suite $(d_{n})$ est \res{géométrique, de raison $0{,}98$}.}{$98\,\%=0{,}98$}
\begin{calculs}
d_{n} &= d_{1}\times q^{n-1} &&\justif{la suite commence à $d_{1}$}\\
d_{n} &= \res{30\times 0{,}98^{n-1}} &&
\end{calculs}
\cas{2)}
\begin{calculs}
d_{5} &= 30\times 0{,}98^{4} &&\justif{$n-1=4$ ; à la calculatrice}\\
d_{5} &\approx \res{27{,}67} &&\justif{arrondi au centième}
\end{calculs}
\explique{Au $5^{\text{e}}$ vol, la batterie tient environ $27{,}67$ minutes.}{}
\cas{3)}
\explique{On cherche le premier $n$ tel que $d_{n}<24$, dans le tableau de valeurs de
$30\times 0{,}98^{x-1}$.}{}
\explique{$d_{12}\approx 24{,}02$ \ et \ $d_{13}\approx 23{,}54<24$.}{on lit les deux
lignes autour du seuil}
\explique{La batterie tient moins de $24$ minutes \res{à partir du $13^{\text{e}}$
vol}.}{$d_{13}$ : le $13^{\text{e}}$ vol, puisque la suite commence à $d_{1}$}
\rem{Erreur fréquente : écrire $d_{n}=30\times 0{,}98^{n}$. La suite commence à $d_{1}$ :
l'exposant est $n-1$.}
\end{exo}

\end{document}
